% concatenated from Joint_Exclusivity_7.tex
\documentclass[letterpaper,12pt]{article}
%\documentclass{statsoc}
%\usepackage{datatool}
\usepackage{diagbox}
\usepackage{latexsym,amssymb,amsmath,bm}
\usepackage{mathtools}
\usepackage{amsthm}
\usepackage{lipsum,graphicx,multicol}
\usepackage{setspace,color}
\usepackage[sort&compress]{natbib}
\usepackage{tcolorbox}
\usepackage{subcaption}
\usepackage{float}
\usepackage[normalem]{ulem}
\usepackage{authblk}
\usepackage{leftidx}
\usepackage{dsfont}
\usepackage{geometry}
\usepackage{bbm}
\usepackage[title,titletoc]{appendix}
\usepackage{actuarialangle}
\usepackage{lscape}
\usepackage{afterpage}
\usepackage{relsize}
\usepackage[colorlinks=true,allcolors=blue, linkcolor=red]{hyperref}
%\usepackage[
%    backend=biber,      % use biber instead of bibtex
%    style=apa,   % or numeric, apa, etc.
%    doi=true,           % show DOIs      
%    natbib=true
%]{biblatex}
%\addbibresource{References.bib}
%geometry{left=2cm,right=2cm,top=2cm,bottom=3cm}
%geometry{left=1.5cm,right=1.5cm,top=2cm,bottom=3cm}
\setlength{\oddsidemargin}{0in}
\setlength{\evensidemargin}{0in}
\setlength{\textheight}{8.5in}
\setlength{\textwidth}{6.7in}
\setlength{\topmargin}{-0.5in}
\allowdisplaybreaks
\renewcommand{\baselinestretch}{1.2}
%\usepackage{stmaryrd}
\usepackage{lineno}
%\linenumbers
\usepackage{bbm}

\usepackage{accents}
\newcommand*{\dt}[1]{%
  \accentset{\mbox{\Large\bfseries .}}{#1}}
\newcommand*{\ddt}[1]{%
  \accentset{\mbox{\large\bfseries .\hspace{-0.25ex}.}}{#1}}
\newcommand\ddfrac[2]{\frac{\displaystyle #1}{\displaystyle #2}}

\def \P {\mathbb{P}}
\def \Re {\mathbb{R}}
\def \ra {\rightarrow}
\def \la {\leftarrow}
\def \oo {\bm{0}}
\def \Chat{{\widehat C}}
\def \Cbar {{\overline C}}
\def \uu {\bm{u}}
\def \T {\mathrm{T}}
\def \tailo {\kappa}
\def \a {\alpha}
\def \rv {\mbox{\rm{rv}}}
\def \al {\alpha}
\def \XX {\bm{X}}
\def \xx {\bm{x}}
\def \ww {\bm{w}}
\def \zz {\bm{z}}
\def \one {\bm{1}}
\def \wbf {\bm{w}}
\def \tailorder {\kappa}
\def \lm {\lambda}
\def \MDA {\mbox{\rm{MDA}}}
\def \eqd {\buildrel d \over =}
\def \rav {\buildrel v \over \rightarrow}

\def \red#1{\textcolor{red}{#1}} % need to solve
\def \blu#1{\textcolor{blue}{#1}} % additions

\def \lmm {\left[}
\def \rmm {\right]}
\def \Fbar {{\overline F}}
\def \Hbar {{\overline H}}
\def \MDA {\mbox{\rm{MDA}}}
\def \pf {\noindent \emph{Proof: }}
\def \QED {\hfill $\square$}
\def \Gbar {{\overline G}}
\def \E {\mathbb{E}}
\def \G {\Gamma}
\def \GUM {\mbox{\rm{Gumbel}}}
\def \FRE {\mbox{\rm{Fr\'{e}chet}}}
\def \lss {\left(}
\def \rss {\right)}
\def \lbb {\left\{}
\def \rbb {\right\}}
\def \half {{\textstyle \frac{1}{2}}}
\def \be {\beta}
\def \d {\textrm{d}}
\def \Ra {\Rightarrow}


\def \aa {\boldsymbol{a}}
\def \bb {\boldsymbol{b}}
\def \cc {\boldsymbol{c}}
\def \dd {\boldsymbol{d}}
\def \ee {\boldsymbol{e}}
\def \ff {\boldsymbol{f}}
\def \gg {\boldsymbol{g}}
\def \hh {\boldsymbol{h}}
\def \ii {\boldsymbol{i}}
\def \jj {\boldsymbol{j}}
\def \kk {\boldsymbol{k}}
\def \ll {\boldsymbol{l}}
\def \mm {\boldsymbol{m}}
\def \nn {\boldsymbol{n}}
\def \ooo {\boldsymbol{o}}
\def \pp {\boldsymbol{p}}
\def \qq {\boldsymbol{q}}
\def \rr {\boldsymbol{r}}
\def \ss {\boldsymbol{s}}
\def \tt {\boldsymbol{t}}
\def \uu {\boldsymbol{u}}
\def \vv {\boldsymbol{v}}
\def \ww {\boldsymbol{w}}
\def \xx {\boldsymbol{x}}
\def \yy {\boldsymbol{y}}
\def \zz {\boldsymbol{z}}
\def \alpp {\boldsymbol{\alpha}}

\def \AA {\boldsymbol{A}}
\def \BB {\boldsymbol{B}}
\def \CC {\boldsymbol{C}}
\def \DD {\boldsymbol{D}}
\def \EE {\boldsymbol{E}}
\def \FF {\boldsymbol{F}}
\def \GG {\boldsymbol{G}}
\def \HH {\boldsymbol{H}}
\def \II {\boldsymbol{I}}
\def \JJ {\boldsymbol{J}}
\def \KK {\boldsymbol{K}}
\def \LL {\boldsymbol{L}}
\def \MM {\boldsymbol{M}}
\def \NN {\mathbb{N}}
\def \OO {\boldsymbol{O}}
\def \PP {\boldsymbol{P}}
\def \QQ {\boldsymbol{Q}}
\def \RR {\mathbb{R}}
\def \SS {\boldsymbol{S}}
\def \TT {\boldsymbol{T}}
\def \UU {\boldsymbol{U}}
\def \VV {\boldsymbol{V}}
\def \WW {\boldsymbol{W}}
\def \XX {\boldsymbol{X}}
\def \YY {\boldsymbol{Y}}
\def \ZZ {\boldsymbol{Z}}
\def\xset{{\mathcal{X}}}
\def\VaR{{\mathrm{VaR}}}
\def\CTE{{\mathrm{CTE}}}
\def\e{{\mathbb E}}
\def\TSD{{\mathrm{TSD}}}
\def\SD{{\mathrm{SD}}}
\def\TCOV{{\mathrm{TCov}}}
\def\Cov{{\mathrm{Cov}}}
\def\Var{{\mathrm{Var}}}
\def\GMD{\mathrm{GMD}}
\def\FOR{\quad \text{for} \quad}

%\def\thetheorem{\arabic{theorem}}
\numberwithin{equation}{section} 
\newtheorem{theorem}{Theorem}[section]
\newtheorem{definition}[theorem]{Definition}
\newtheorem{question}{Question}
\newtheorem{example}[theorem]{Example}
\newtheorem{remark}[theorem]{Remark}
\newtheorem{proposition}[theorem]{Proposition}
\newtheorem{lemma}[theorem]{Lemma}
\newtheorem{corollary}[theorem]{Corollary}
\newtheorem{property}{Property}
%\setcounter{theorem}{0}
\newcommand{\N}{{N}}
\newcommand{\M}{{M}}
\newcommand{\J}{{J}}
\newcommand{\K}{{K}}
%\newcommand{\Mi}[1]{\mathcal{M}^{\mkern 1mu#1}}
%\newcommand{\Mj}[1]{\mathcal{M}_{#1}}
%\newcommand{\Mij}[2]{\mathcal{M}^{\mkern 1mu#1}_{#2}}
\newcommand{\I}{{I}}
\newcommand{\F}{\,\overline{\kern-0.175em F\kern-0em}}
\newcommand{\xb}{\,\overline{\kern-0.125em x\kern-0em}}
\newcommand{\Xb}{\,\overline{\kern-0.125em X\kern-0em}}


\makeatletter
\setlength{\arraycolsep}{2\p@} % make spaces around "=" in eqnarray smaller
\makeatother

%\newcommand{\JS}[1]{\textcolor{orange}{#1}}
\definecolor{darkread}{rgb}{0.7, 0, 0}
\newcommand{\EF}[1]{\textcolor{darkread}{#1}}
\newcommand{\AF}[1]{\textcolor{green}{#1}}


%\newcommand{\rd}[1]{\textcolor{darkread}{#1}}

\begin{document}
\doublespace
\title{\Large\textbf{{Joint Exclusivity}}}

\author[]{\small Nawaf Mohammed \thanks{\url{nawaf.mohammed.ac@gmail.com}}}

\affil[]{\footnotesize }
\date{}
\maketitle
\vspace{-1cm}
\begin{abstract}
We introduce \emph{joint exclusivity} (JE), a flexible extension of mutual exclusivity for non-negative random vectors that prohibits only the simultaneous positivity of all $n$ components. Unlike mutual exclusivity, JE admits a broad class of marginals while retaining a sharp existence criterion. We prove that a JE random vector with prescribed marginals $F_1,\ldots,F_n$ exists if and only if
\[
\sum_{i=1}^n \F_i(0)\le n-1.
\]
We also provide a canonical face-based construction with freely specified copulas within each face, and show that the trivariate case can elegantly be reduced to a one-parameter family. Finally, we extend the framework through Wang-type distortions and an $m$-exclusivity hierarchy interpolating between mutual exclusivity and JE.


\vspace{5cm}
\noindent\textit{Key words and phrases}: counter-monotonicity; joint exclusivity; mutual exclusivity; $m$-exclusivity; generalised joint exclusivity; joint mixability; distortion functions.
\end{abstract}

\newpage
{\color{black}
\section{Introduction}
\label{sec:intro}

Dependence modelling plays a central role in quantitative risk analysis, where the choice of joint distribution
profoundly shapes aggregate loss behaviour and the values of associated risk
measures~\citep{Lauzier2023, Deelstra2011, Dhaene2002a, Cheung2014a}. Among
all dependence structures, those of an extremal nature are of special importance:
they yield sharp bounds on risk measures and serve as natural benchmarks for
sensitivity analysis.

In the bivariate setting, the positive extreme is \emph{co-monotonicity}
\citep{Dhaene2002, Nelson2010}, which captures the notion that both components
always move in the same direction. Formally, $X_{12}=(X_1, X_2)$ is co-monotonic
if
\[
    X_{12} \stackrel{d}{=} \bigl(h_1(Z),\, h_2(Z)\bigr),
\]
for some random variable $Z$ and functions $h_1$, $h_2$ that are both
non-decreasing (or both non-increasing). This structure is characterised by the
Fr\'{e}chet upper bound
\[
    F_{12}(x_1, x_2) = \min\left\{F_1(x_1),\, F_2(x_2)\right\},
\]
where $F_{12}$ denotes the joint cumulative distribution function (CDF) of
$X_{12}$ and $F_1$, $F_2$ are the marginal distributions. Co-monotonicity
represents the strongest form of positive dependence: an increase in one
component almost surely does not decrease the other.

At the opposite extreme lies \emph{counter-monotonicity} \citep{Dhaene1999,
Nelson2010}, which captures perfect negative association. A bivariate vector
$X_{12}$ is counter-monotonic if
\[
    X_{12} \stackrel{d}{=} \bigl(h_1(Z),\, h_2(Z)\bigr),
\]
where $h_1$ is non-decreasing and $h_2$ is non-increasing (or vice versa), for
some random variable $Z$. Its joint distribution is given by the Fr\'{e}chet
lower bound
\[
    F_{12}(x_1, x_2) = \max\left\{F_1(x_1) + F_2(x_2) - 1,\; 0\right\}.
\]
In the bivariate case, both Fr\'{e}chet bounds define valid joint CDFs for
arbitrary marginals $F_1$ and $F_2$, so the two extremal dependence structures
are uniquely and unambiguously defined.

These two notions, however, behave very differently in higher dimensions. Let
$\N = \{1,\dots, n\}$ denote the index set. Co-monotonicity extends naturally to
the $n$-variate setting: for $X_{\N}= (X_1, \dots, X_n)$, the representation
\[
    X_{\N}\stackrel{d}{=} \bigl(h_1(Z),\, \dots,\, h_n(Z)\bigr),
\]
where all $h_i$ are non-decreasing (or all non-increasing), remains well-defined
for any $n \ge 2$. Moreover, the Fr\'{e}chet upper bound
\[
    F_{\N}(x_1, \dots, x_n) = \min_{i\in\N} F_i(x_i)
\]
continues to define a valid joint distribution for arbitrary marginals and any
$n \ge 2$.

Counter-monotonicity, by contrast, does not admit a canonical multivariate
extension. For $n \ge 3$, if two pairs
of components are each driven by opposing monotone functions of a common variable,
the induced relationship between the remaining pair is necessarily co-monotone,
yielding intricacies that are handled non-uniquely. Correspondingly, the multivariate
Fr\'{e}chet lower bound
\[
    F_{\N}(x_1, \dots, x_n)
        = \max\left\{\sum_{i\in\N} F_i(x_i) - (n-1),\; 0\right\}
\]
fails, in general, to define a valid joint distribution when $n \ge 3$. As a
result, several extensions of counter-monotonicity have been proposed in the
literature, each tailored to specific modelling contexts.

Among the strongest of these is \emph{mutual exclusivity} (ME) \citep{Dhaene1999,
Cheung2014}. Letting $\F$ denote the decumulative distribution function (DDF),
a non-negative random vector $X_{\N}$ is ME if
\begin{equation}
    \label{eq:MEdef}
    \F_{ij}(0,0) = 0, \qquad \text{for all } i<j,\quad i,j \in \N,
\end{equation}
that is, no two components are simultaneously positive. An ME random vector with
prescribed marginals $F_1, \dots, F_n$ exists if and only if
\begin{equation}
    \label{eq:MEcondition}
    \sum_{i\in\N}\F_i(0) \le 1.
\end{equation}
Whenever \eqref{eq:MEcondition} holds, the Fr\'{e}chet lower bound is a valid
joint distribution and coincides with that of the ME random vector. This
tractability is one of the principal appeals of ME: existence is fully determined
by the single inequality~\eqref{eq:MEcondition}, and the joint distribution is
directly given by the Fr\'{e}chet lower bound. However,
condition~\eqref{eq:MEcondition} is rather restrictive---it requires that the
aggregate probability of any component being strictly positive does not exceed
unity---which considerably limits the practical applicability of ME.

Alternative notions of counter-monotonicity have been introduced to overcome this
limitation, including joint mixability \citep{Wang2016}, $d$-counter-monotonicity
\citep{Lee2014}, and $\Sigma$-counter-monotonicity; see \cite{Puccetti2015} for a
unified treatment. These approaches typically impose that a prescribed function of
the components is almost surely constant, thereby permitting a broader class of
marginals. However, they are often difficult to construct explicitly or to
characterise completely for a given set of marginals. This stands in sharp
contrast to ME, whose existence is fully determined by condition~\eqref{eq:MEcondition}
and whose joint distribution is directly given by the Fr\'{e}chet lower bound.
This observation motivates the following question: \emph{does there exist a notion
of counter-monotonicity that is more flexible than ME, yet remains amenable to
explicit construction and admits a complete, tractable characterisation analogous
to~\eqref{eq:MEcondition}?}

A natural answer emerges by relaxing the pairwise requirement
in~\eqref{eq:MEdef}. Rather than prohibiting any pair of components from being
simultaneously positive, one may instead require only that all $n$ components are
never simultaneously positive. This leads to the following definition.

\begin{definition}
\label{def:JE}
A non-negative random vector $X_{\N}$ is said to be \emph{jointly exclusive} (JE)
if
\begin{equation}
    \label{eq:JEdef}
    \F_{\N}(0,\dots,0) = 0,
    \text{ or equivalently, }
    \prod_{i\in\N}X_i=0\text{ almost surely},
\end{equation}
that is, its support is contained in $[0,\infty)^n \setminus (0,\infty)^n$.
\end{definition}

\begin{figure}[htbp]
\centering
\includegraphics[width=\textwidth]{MEvsJE.pdf}
\caption{Comparison of the supports of ME and JE random vectors for $n=3$.}
\label{fig:MEvsJE}
\end{figure}

When $n = 2$, conditions \eqref{eq:MEdef} and \eqref{eq:JEdef} coincide, so ME
and JE are equivalent. For $n \ge 3$, ME implies JE but not conversely, making JE
a strictly more flexible notion; see Figure~\ref{fig:MEvsJE} for an illustration
of the distinction between the two supports when $n = 3$. As we demonstrate in
this paper, JE additionally admits a canonical construction and a complete
characterisation of the admissible marginals, thereby providing an affirmative
answer to the question above.

The remainder of the paper is organised as follows. Section~\ref{sec:preliminaries}
introduces the notation and presents basic properties of JE.
Section~\ref{sec:main_results} develops the main construction methods and
characterisation results. Further remarks on the reflected and translated variants
of JE as well as its connection to joint mixability are discussed in
Section~\ref{sec:further_discussion}. Section~\ref{sec:conclusions} concludes.
\section{Preliminaries}
\label{sec:preliminaries}

Throughout the paper we work with the index set $\N=\{1,\dots,n\}$. Define
$\M=\{\I\subset \N : 2\le |\I|\le n-1\}$ as the collection of all proper subsets
of $\N$ of intermediate cardinality, and for each $i\in\N$ let
$\M_i=\{\I\in \M : i\in \I\}$ denote those containing $i$; the latter family
arises naturally in the constructions of Section~\ref{sec:main_results}. We
primarily consider non-negative random vectors $X_{\N}=(X_1,\dots,X_n)$, with
sub-vectors $X_{\I}=(X_i,\,i\in\I)$ for $\I\subseteq \N$, and write
$x_{\N}=(x_1,\dots,x_n)\in\RR^n$ and $x_{\I}=(x_i,\,i\in\I)$ for their
deterministic counterparts. We use $F$, $\F$, $f$, and $\varphi$ to denote the
CDF, DDF, density function, and characteristic function, respectively, with
subscripts identifying the corresponding vector or sub-vector; for example,
$F_{\N}$ and $\F_{\N}$ denote the CDF and DDF of $X_{\N}$, and analogously for
$F_{\I}$ and $\F_{\I}$.

By Definition~\ref{def:JE} and as illustrated in Figure~\ref{fig:MEvsJE}, the
support of a JE random vector is contained in
$\bigcup_{i\in\N}\{x\in[0,\infty)^n : x_i=0\}$, a union of $n$ coordinate
hyperplanes of dimension $n-1$ within $[0,\infty)^n$. The joint distribution
therefore assigns no $n$-dimensional Lebesgue mass and is entirely characterised
by its lower-dimensional marginals, as formalised below.

\begin{proposition}
\label{prop:JECDF}
A non-negative random vector $X_{\N}$ is JE if and only if its CDF satisfies
\begin{equation}
\label{eq:JECDF}
F_{\N}(x_{\N})=\begin{cases}
1-\sum\limits_{i\in\N}\F_i(x_i)+\sum\limits_{\I\in\M}(-1)^{|\I|}\F_{\I}(x_{\I})\\[4pt]
\hspace*{0.65cm}=(-1)^{n-1}\left(1-\sum\limits_{i\in\N}F_i(x_i)
+\sum\limits_{\I\in\M}(-1)^{|\I|}F_{\I}(x_{\I})\right), & x_{\N}\in[0,\infty)^n, \\[8pt]
0, & \text{otherwise,}
\end{cases}
\end{equation}
and if and only if its characteristic function satisfies
\begin{equation}
\label{eq:JECF}
\varphi_{\N}(t_{\N})
= (-1)^{n-1}\left(1-\sum_{i\in\N}\varphi_i(t_i)
+\sum_{\I\in\M}(-1)^{|\I|}\,\varphi_{\I}(t_{\I})\right),
\quad t_{\N}\in\RR^n.
\end{equation}
\end{proposition}

\begin{proof}
\textit{CDF representation.}
For $x_{\N}\in[0,\infty)^n$, the inclusion--exclusion principle yields
\begin{equation}
\label{eq:IEidentityDDF}
F_{\N}(x_{\N}) = \sum_{J\subseteq\N}(-1)^{|J|}\F_J(x_J), \qquad \F_\emptyset\equiv 1.
\end{equation}
Under JE, non-negativity gives $\F_{\N}(x_{\N})=\P(X_i>x_i,\,\forall
i\in\N)\le \P\bigl(\prod_{i\in\N}X_i>0\bigr)=0$ for all $x_\N\in[0,\infty)^n$.
Dropping the $J=\N$ term from~\eqref{eq:IEidentityDDF} yields the first expression
in~\eqref{eq:JECDF}. Conversely, if $F_\N$ satisfies this expression, note that
it equals $\sum_{J\subset\N}(-1)^{|J|}\F_J(x_J)$; substituting into the
right-hand side of~\eqref{eq:IEidentityDDF} then gives $(-1)^n\F_\N(x_\N)=0$,
confirming JE. For $x_\N\notin[0,\infty)^n$, $F_\N(x_\N)=0$ follows immediately
from non-negativity.

For the second expression, the dual inclusion--exclusion identity
$\F_{\N}(x_\N)=\sum_{J\subseteq\N}(-1)^{|J|}F_J(x_J)$, combined with
$\F_\N(x_\N)=0$, gives $(-1)^nF_\N(x_\N)=-\sum_{J\subset\N}(-1)^{|J|}F_J(x_J)$,
from which
$F_\N(x_\N)=(-1)^{n-1}\sum_{J\subset\N}(-1)^{|J|}F_J(x_J)$,
yielding the second expression.

\textit{Characteristic function representation.}
Under JE, $\prod_{i\in\N}X_i=0$ a.s., so for any $t_\N\in\RR^n$ at least one
factor satisfies $\mathrm{e}^{\mathrm{i}t_i\cdot 0}-1=0$, giving
$\prod_{i\in\N}(\mathrm{e}^{\mathrm{i}t_iX_i}-1)=0$ a.s.
Expanding via the identity
$\prod_{i\in\N}(c_i-1)=\sum_{J\subseteq\N}(-1)^{n-|J|}\prod_{j\in J}c_j$
and taking expectations yields
$\sum_{J\subseteq\N}(-1)^{n-|J|}\varphi_J(t_J)=0$.
Isolating the $J=\N$ term and rearranging gives~\eqref{eq:JECF}. The converse
follows from the uniqueness theorem for characteristic functions together with the
CDF equivalence established above.
\end{proof}

\begin{remark}
\label{rm:JECDFrm}
The representation~\eqref{eq:JECDF} illuminates the structural distinction between
JE and ME. Under ME, no two components are simultaneously positive, so
$\F_{\I}(x_\I)=0$ for every $\I\in\M$ and $x_\I\in[0,\infty)^{|\I|}$. The
interaction terms
\begin{equation}
\label{eq:MECDF_Example}
\sum_{\I\in\M}(-1)^{|\I|}\F_{\I}(x_{\I})
\end{equation}
therefore vanish identically, and the CDF reduces to
\[
F_{\N}(x_{\N})=\begin{cases}
1-\sum\limits_{i\in\N}\F_i(x_i)
=\sum\limits_{i\in\N}F_i(x_i)-(n-1), & x_{\N}\in[0,\infty)^n,\\[4pt]
0, & \text{otherwise,}
\end{cases}
\]
which coincides with the Fr\'{e}chet lower bound and defines a valid CDF under
condition~\eqref{eq:MEcondition}. Under JE, by contrast, the interaction
terms~\eqref{eq:MECDF_Example} need not vanish, reflecting the richer dependence
structure that permits components to be simultaneously positive in pairs or smaller
sub-groups while still satisfying the global JE constraint.
\end{remark}
%+++++++++++++++++++++++++++++++++++++++++++++++++++
\section{Main Results}
\label{sec:main_results}

We begin by establishing the fundamental characterisation of JE in terms of the
admissible marginal distributions.

\begin{theorem}
\label{thm:JE_Characterisation_1}
Let $F_1,\dots,F_n$ be marginal distributions on $[0,\infty)$. A JE random
vector $X_{\N}$ with marginals $F_1,\dots,F_n$ exists if and only if
\begin{equation}
\label{eq:JECondition}
\sum_{i\in\N}\F_i(0)\le n-1.
\end{equation}
\end{theorem}

\begin{proof}
\textit{Necessity.}
Suppose that $X_{\N}$ is JE with marginals $F_1,\dots,F_n$. By
Definition~\ref{def:JE}, $\F_{\N}(0,\dots,0)=0$. The Fr\'{e}chet lower bound
for the joint DDF gives
\[
0 \le \max\left\{\sum_{i\in\N}\F_i(0)-(n-1),\;0\right\}
    \le \F_{\N}(0,\dots,0)=0,
\]
from which $\sum_{i\in\N}\F_i(0)\le n-1$ follows immediately.

\medskip
\noindent\textit{Sufficiency.}
Assume that~\eqref{eq:JECondition} holds. We construct a valid JE distribution
with marginals $F_1,\dots,F_n$ by assigning probability mass to each piece of the
partition
\[
[0,\infty)^n\setminus(0,\infty)^n
= \{0_{\N}\}
  \cup\bigcup_{i\in\N} A_i
  \cup\bigcup_{\I\in\M} E_{\I},
\]
where $0_{\N}$ denotes the origin, $A_i=\{x_{\N}:x_i>0,\,x_j=0\ \forall j\ne i\}$
is the $i$-th open semi-axis, and
\[
E_{\I}=\{x_{\N}: x_i>0\ \forall i\in\I,\;x_j=0\ \forall j\in\N\setminus\I\}
\]
is the open face indexed by $\I\in\M$. These sets are mutually disjoint.

\smallskip
\noindent\textit{Step~1: Intermediate faces.}
For each $\I\in\M$ and $x_{\I}\in[0,\infty)^{|\I|}$, set
\begin{equation}
\label{eq:JEProbFaces}
\P\left(X_i>x_i,\;i\in \I;\;X_j=0,\;j\in\N\setminus \I\right)
= p_{\I}\,C_{\I}\left(\frac{\F_i(x_i)}{\F_i(0)},\;i\in\I\right),
\end{equation}
where $p_{\I}\ge 0$ is the total probability mass assigned to face $E_\I$ and
$C_{\I}$ is an arbitrary copula on $[0,1]^{|\I|}$ \citep{Nelson2010}. The
argument $\frac{\F_i(x_i)}{\F_i(0)}\in[0,1]$ is the survival function of $X_i$
conditional on $X_i>0$; when $\F_i(0)=0$ for some $i\in\I$, the
constraint~\eqref{eq:JE_parameters_1} forces $p_\I=0$, so the formula need only
be evaluated when all denominators are strictly positive. Since the
faces $\{E_{\I}:\I\in\M\}$ are disjoint, the copulas $\{C_\I:\I\in\M\}$ may be
chosen independently; in particular, $C_{\I_1}$ need not be a marginal of
$C_{\I_2}$ even when $\I_1\subset\I_2$.

\smallskip
\noindent\textit{Step~2: Coordinate axes.}
The marginal constraint $\P(X_i>x_i)=\F_i(x_i)$ must hold for each $i\in\N$.
The contribution of face $E_{\I}$ (for $\I\ni i$) to $\P(X_i>x_i)$ is obtained by
setting $x_k=0$ (and hence $\frac{\F_k(x_k)}{\F_k(0)}=1$) for all
$k\in\I\setminus\{i\}$ in~\eqref{eq:JEProbFaces}, which by the copula
marginalisation property $C(u,1,\dots,1)=u$ yields $p_\I\cdot\frac{\F_i(x_i)}{\F_i(0)}$.
Summing these face contributions over all $\I\in\M_i$ and requiring equality with
$\F_i(x_i)$ uniquely determines the axis mass:
\begin{equation}
\label{eq:JEProbAxes}
\P\left(X_i>x_i;\;X_j=0,\;j\in\N\setminus\{i\}\right)
=\F_i(x_i)\left(1-\frac{\displaystyle\sum_{\I\in \M_i}p_{\I}}{\F_i(0)}\right),
\quad x_i\in[0,\infty).
\end{equation}
This expression defines a valid (non-negative, non-increasing) survival function
for the $i$-th axis component if and only if
\begin{equation}
\label{eq:JE_parameters_1}
\sum_{\I\in \M_i}p_{\I}\le\, \F_i(0), \qquad i\in\N.
\end{equation}

\smallskip
\noindent\textit{Step~3: Origin.}
The total probability mass on all axes equals
$\sum_{i\in\N}\bigl[\F_i(0)-\sum_{\I\in\M_i}p_\I\bigr]
=\sum_{i\in\N}\F_i(0)-\sum_{\I\in\M}|\I|\,p_\I$,
and the total mass on all faces equals $\sum_{\I\in\M}p_\I$.
Subtracting from unity, the origin receives mass
\begin{equation}
\label{eq:JEProbOrigin}
\P\left(X_i=0,\;i\in\N\right)
=1-\sum_{i\in\N}\F_i(0)+\sum_{\I\in\M}(|\I|-1)\,p_{\I},
\end{equation}
which is non-negative if and only if
\begin{equation}
\label{eq:JE_parameters_2}
\sum_{\I\in\M}(|\I|-1)\,p_{\I}\ge \sum_{i\in\N}\F_i(0)-1.
\end{equation}
One can verify that the total probability sums to unity under this
assignment. The construction therefore yields a valid JE distribution if and
only if there exist coefficients $p_\I\ge 0$ satisfying both
\eqref{eq:JE_parameters_1} and~\eqref{eq:JE_parameters_2}.

\smallskip
\noindent\textit{Step~4: Feasibility via linear programming.}
The question reduces to whether the maximum of $\sum_{\I\in\M}(|\I|-1)p_\I$
over the feasible set defined by~\eqref{eq:JE_parameters_1} is at least
$\sum_{i\in\N}\F_i(0)-1$. Consider the primal linear program
\begin{gather}
\label{sys:original}
\text{maximise}\quad\sum_{\I\in\M}(|\I|-1)p_{\I}, \\
\nonumber
\text{subject to}\quad p_{\I}\ge 0,\quad
\sum_{\I\in \M_i}p_{\I}\le\, \F_i(0),\quad i\in\N,
\end{gather}
with dual problem
\begin{gather}
\label{sys:dual}
\text{minimise}\quad\sum_{i\in\N}\F_i(0)\,r_i, \\
\nonumber
\text{subject to}\quad r_i\ge 0,\quad
\sum_{i\in \I}r_i\ge |\I|-1,\quad \I\in\M.
\end{gather}
Both programs are feasible: $p_\I=0$ is primal feasible, and setting $r_i=1$
for all $i$ is dual feasible, since $|\I|\ge|\I|-1$ for every $\I\in\M$. By
strong duality, the optimal values of~\eqref{sys:original} and~\eqref{sys:dual}
coincide, so the sufficiency condition becomes
\[
\min_{r\text{ dual-feasible}}\sum_{i\in\N}\F_i(0)\,r_i\ge\sum_{i\in\N}\F_i(0)-1.
\]

The dual minimum is attained at a vertex of the feasible polytope. An analysis
of the extreme points of the dual polytope reveals that the minimum of the dual
objective is achieved either at the symmetric point
$r^*=\left(\frac{n-2}{n-1},\dots,\frac{n-2}{n-1}\right)$
or at one of the $n$ permutations of $(1,\dots,1,0)$, with the zero in
position~$j$. Feasibility of both candidates is immediate: at $r^*$,
$\sum_{i\in\I}r_i^*=|\I|\cdot\frac{n-2}{n-1}\ge|\I|-1$ holds for all
$\I\in\M$ because $|\I|\le n-1$; at $(1,\dots,1,0_j,1,\dots,1)$,
the dual constraint evaluates to $|\I|$ or $|\I|-1$ according to whether
$j\notin\I$ or $j\in\I$, both of which satisfy $\ge|\I|-1$. The corresponding
objective values are $\frac{n-2}{n-1}\sum_{i\in\N}\F_i(0)$ and
$\sum_{i\in\N}\F_i(0)-\F_j(0)$; minimising the latter over $j$ gives
$\sum_{i\in\N}\F_i(0)-\max_{i\in\N}\F_i(0)$. Hence
\begin{equation}
\label{eq:JE_minimum_sol}
\min_{r\text{ dual-feasible}}\sum_{i\in\N}\F_i(0)\,r_i
=\min\left\{\frac{n-2}{n-1}\sum_{i\in\N}\F_i(0),\;
\sum_{i\in\N}\F_i(0)-\max_{i\in\N}\F_i(0)\right\}.
\end{equation}

It remains to verify that this minimum is at least $\sum_{i\in\N}\F_i(0)-1$ as \eqref{eq:JE_parameters_2} dictates.
If the minimum in~\eqref{eq:JE_minimum_sol} is attained by the second term, the
requirement reduces to $\max_{i\in\N}\F_i(0)\le 1$, which holds trivially since
each $\F_i(0)$ is a probability. If it is attained by the first term, the
requirement becomes
\[
\frac{n-2}{n-1}\sum_{i\in\N}\F_i(0)\ge\sum_{i\in\N}\F_i(0)-1,
\]
which simplifies to $\sum_{i\in\N}\F_i(0)\le n-1$, precisely
condition~\eqref{eq:JECondition}. In both cases, \eqref{eq:JECondition} is
sufficient to guarantee the existence of feasible coefficients $p_\I$, and the
construction in Steps~1--3 defines a valid JE random vector with the prescribed
marginals $F_1,\dots,F_n$.
\end{proof}


%+++++++++++++++++++++++++++++++++++++++++++++++++++
The necessary and sufficient condition~\eqref{eq:JECondition} for JE is
significantly less restrictive than its ME counterpart~\eqref{eq:MEcondition},
and the gap widens with increasing dimension $n$. Equivalently in terms of
the marginal CDFs, the JE condition reads
\[
\sum_{i\in\N} F_i(0) \ge 1,
\]
while ME requires the substantially stronger
\[
\sum_{i\in\N} F_i(0) \ge n-1.
\]
Under ME, the constraint $\sum_{i\in\N}\F_i(0)\le 1$ precludes any marginal from
being free of an atom at zero unless all other components are almost surely zero:
if $\F_{i_0}(0)=1$ for some $i_0\in\N$, then $\sum_{j\ne i_0}\F_j(0)\le0$,
forcing $X_j=0$ almost surely for every $j\ne i_0$. Hence, for non-degenerate
marginals, ME requires $F_i(0)>0$ for every $i\in\N$. Under JE, by contrast,
few marginals need carry atoms at zero. If $n-2$ marginals are fully supported on
$(0,\infty)$, each contributes $\F_i(0)=1$, so condition~\eqref{eq:JECondition}
reduces to
\[
\F_{i_1}(0)+\F_{i_2}(0)\le 1,
\quad\text{equivalently}\quad
F_{i_1}(0)+F_{i_2}(0)\ge 1,
\]
for the two remaining marginals. Hence, as few as two marginals need carry atoms
at zero, subject to the mild joint condition that their point masses at zero sum
to at least one, while the remaining $n-2$ marginals may be fully supported on
$(0,\infty)$.

\medskip

To further investigate the structural properties of JE, we adopt the construction
introduced in the proof of Theorem~\ref{thm:JE_Characterisation_1} as a canonical
model.

\begin{definition}
\label{def:JECanonical}
A random vector $X_{\N}$ is called a \emph{canonical JE random vector} if its
distribution is specified by \eqref{eq:JEProbFaces}, \eqref{eq:JEProbAxes}, and \eqref{eq:JEProbOrigin} for some coefficients $p_{\I}\in[0,1]$, $\I\in\M$, satisfying
\eqref{eq:JE_parameters_1} and~\eqref{eq:JE_parameters_2}.
\end{definition}

For fixed marginals $F_1,\dots,F_n$, if condition~\eqref{eq:JECondition} holds,
Theorem~\ref{thm:JE_Characterisation_1} guarantees the existence of a canonical JE
random vector. Its CDF admits the following explicit form, as a direct consequence
of Proposition~\ref{prop:JECDF}.

\begin{corollary}
\label{cor:JE_Canonical_CDF_phi}
A random vector $X_{\N}$ is canonical JE if and only if its joint CDF is given by
\begin{equation}
\label{eq:JE_Canonical_CDF}
F_{\N}(x_{\N})=\begin{cases}
{\displaystyle
1-\sum_{i\in\N}\F_i(x_i)
+\sum_{\I\in\M}(-1)^{|\I|}
\sum_{\substack{\J\in\M\\ \I\subseteq \J}}
p_{\J}\,C_{\J}\left(\frac{\F_i(x_i)}{\F_i(0)},\,i\in\I;\;1,\,j\in \J\setminus \I\right)},
\\[10pt]
\hspace{10.5cm} x_{\N}\in[0,\infty)^n, \\[6pt]
0, \hspace{9.45cm} \text{otherwise.}
\end{cases}
\end{equation}
\end{corollary}

\begin{proof}
For any $\I\in\M$ and $x_\I\in[0,\infty)^{|\I|}$, the canonical construction
concentrates mass for the event $\{X_i>x_i,\,i\in\I\}$ exclusively on faces
$G_\J$ with $\J\supseteq\I$, so that
\[
\F_\I(x_\I)
=\sum_{\substack{\J\in\M\\\J\supseteq\I}}
p_\J\,C_\J\left(\frac{\F_i(x_i)}{\F_i(0)},\,i\in\I;\;1,\,j\in\J\setminus\I\right),
\]
where the arguments equal one for $j\in\J\setminus\I$ because $x_j=0$ implies
$\F_j(x_j)/\F_j(0)=1$, by the copula marginalisation property
$C(u,1,\dots,1)=u$. Substituting this expression into the first line
of~\eqref{eq:JECDF} yields~\eqref{eq:JE_Canonical_CDF}. Conversely, given a CDF
of the form~\eqref{eq:JE_Canonical_CDF}, the face probabilities $\{p_\J\}$ and
copulas $\{C_\J\}$ are recoverable from the joint DDFs $\F_\I$, confirming that
the random vector is canonical JE.
\end{proof}

For canonical JE random vectors, the coefficients $p_{\I}$ determine how
probability mass is allocated across the faces of the support, and thereby
indirectly along the axes and at the origin. The admissible values of
$\{p_\I\}_{\I\in\M}$ are governed entirely by the marginal
quantities $\F_1(0),\dots,\F_n(0)$, as encoded in
conditions~\eqref{eq:JE_parameters_1} and~\eqref{eq:JE_parameters_2}. In contrast,
the copulas $C_{\I}$ control the internal dependence structure within each face
$E_{\I}$; they may be chosen independently across different $\I\in\M$ and are not
constrained by the marginals.

The following remark illustrates several instructive choices of the parameters
$p_{\I}$ and their implications for the resulting dependence structure.

\begin{remark}
\label{rm:JE_Canonical_Remark}
\mbox{}
\begin{itemize}

\item[(1)]
If $p_{\I}=0$ for some $\I\in\M$, then face $E_{\I}$ is excluded from the support.
Moreover, the Fr\'{e}chet lower bound applied to~\eqref{eq:JEProbFaces} at
$x_\I=0$ gives
\[
\max\left\{
\sum_{i\in\I}\F_i(0)+\sum_{j\in\N\setminus \I}\bigl(1-\F_j(0)\bigr)-(n-1),\;0
\right\}\le p_{\I}=0,
\]
which yields the necessary condition
\[
\sum_{i\in\I}\F_i(0)-\sum_{j\in\N\setminus \I}\F_j(0)\le |\I|-1.
\]
At the extreme, if $p_{\I}=0$ for all $\I\in\M$, all probability mass is
concentrated on the axes and the origin, and no two components are simultaneously
positive. This recovers the ME structure, and condition~\eqref{eq:JE_parameters_2}
reduces to
\[
\sum_{i\in\N}\F_i(0)\le 1,
\]
which is precisely the ME condition~\eqref{eq:MEcondition}.

\item[(2)]
If $p_{\I}=1$ for some $\I\in\M$, all probability mass is concentrated on face
$E_{\I}$, and necessarily $p_{\J}=0$ for all $\J\in\M\setminus\{\I\}$. The
Fr\'{e}chet upper bound applied to~\eqref{eq:JEProbFaces} at $x_\I=0$ requires
\[
1=p_{\I}\le
\min\left\{\min_{i\in\I}\F_i(0),\;\min_{j\in\N\setminus\I}\bigl(1-\F_j(0)\bigr)\right\},
\]
which implies $\F_i(0)=1$ (i.e.\ $\P(X_i>0)=1$) for all $i\in\I$, and
$\F_j(0)=0$ (i.e.\ $X_j=0$ almost surely) for all $j\in\N\setminus\I$.

\item[(3)]
The mass on the $i$-th axis is given by~\eqref{eq:JEProbAxes}. Zero axis mass
for component $i$ requires equality in~\eqref{eq:JE_parameters_1}:
$\sum_{\I\in \M_i}p_{\I}=\F_i(0)$. Imposing this simultaneously for all
$i\in\N$ yields a system of $n$ linear equations in $|\M|=2^n-n-2$ unknowns.
For $n\ge 4$, the system is under-determined and generally admits infinitely many
solutions.

For $n=3$, the system is square ($|\M|=3$ equations in 3 unknowns) and has the
unique solution
\[
p_{\I}=\frac{\sum_{i\in\I}\F_i(0)-\F_j(0)}{2},
\qquad j\in\N\setminus\I,\quad\I\in\M.
\]
This solution satisfies $p_\I\ge0$ for all $\I\in\M$ if and only if
$\max_{k\in\N}\F_k(0)\le\sum_{j\ne k}\F_j(0)$ for every $k\in\N$, that is,
the three quantities $\F_1(0),\F_2(0),\F_3(0)$ satisfy the triangle inequality.
When this condition fails, simultaneous zero axis mass across all three axes is
infeasible.

\item[(4)]
The mass at the origin is governed by~\eqref{eq:JEProbOrigin}. Provided that
\eqref{eq:JE_parameters_1} holds, assigning zero mass to the origin requires
equality in~\eqref{eq:JE_parameters_2}:
\[
\sum_{\I\in\M}(|\I|-1)\,p_{\I}= \sum_{i\in\N}\F_i(0)-1.
\]

\item[(5)]
Suppose $\sum_{i\in\N}\F_i(0)=n-1$, so that equality is achieved
in~\eqref{eq:JECondition}. Let $c_0=\P(X_i=0,\,i\in\N)$ denote the origin mass
and $c_i=\P(X_i>0;\,X_j=0,\,j\in\N\setminus\{i\})$ the mass on the $i$-th axis.
Summing~\eqref{eq:JEProbAxes} over all $i\in\N$ and evaluating at $x_i=0$ gives
\begin{equation}
\label{eq:remark5_identity}
\sum_{i\in\N}c_i+\sum_{\I\in \M}|\I|\,p_{\I}=n-1.
\end{equation}
Since $|\I|\le n-1$ for every $\I\in\M$, and the total face mass satisfies
$\sum_{\I\in\M}p_\I=1-c_0-\sum_{i\in\N}c_i$ by normalization, we have
\[
\sum_{\I\in\M}|\I|\,p_\I\le (n-1)\sum_{\I\in\M}p_\I
=(n-1)\left(1-c_0-\sum_{i\in\N}c_i\right).
\]
Substituting into~\eqref{eq:remark5_identity} and rearranging yields
\[
(n-1)\,c_0+(n-2)\sum_{i\in\N}c_i\le 0.
\]
Since $n\ge 2$ implies $n-1\ge1$ and $n-2\ge0$, and all terms are
non-negative, it follows that $c_0=0$ and $c_i=0$ for every $i\in\N$. Consequently,
when equality holds in~\eqref{eq:JECondition}, the JE random vector concentrates
its mass entirely on the faces, with $\sum_{\I\in\M}p_\I=1$.
\end{itemize}
\end{remark}

Once the marginals are fixed, the admissible coefficients $\{p_\I\}_{\I\in\M}$ are
precisely those satisfying~\eqref{eq:JE_parameters_1} and~\eqref{eq:JE_parameters_2}.
In general, these constraints define a polytope in $\RR^{|\M|}$, leaving
substantial freedom in the choice of dependence structure. For $n=3$, however,
this polytope has a particularly tractable structure. As noted in
Remark~\ref{rm:JE_Canonical_Remark}-(3), the sum in~\eqref{eq:JE_parameters_2}
equals one half of the aggregate of the three sums in~\eqref{eq:JE_parameters_1},
rendering one constraint redundant. The system then reduces to three independent
linear constraints in three unknowns, and the feasible set collapses from a
three-dimensional region to a one-dimensional affine family. In particular, all
three coefficients $p_\I$ can be jointly parametrised by a single scalar
$\lambda\in[0,1]$, interpolating linearly between their extremal values implied
by~\eqref{eq:JE_parameters_1} and~\eqref{eq:JE_parameters_2}. The following
proposition makes this structure explicit.

\begin{proposition}
\label{prop:JE_trivariate}
Let $n=3$ and suppose that condition~\eqref{eq:JECondition} holds. For each
$\I\in\M$, define
\begin{equation}
\label{eq:JE_trivariate_P}
p_{\I}(\lambda)
=\lambda\,\min_{i\in\I}\F_i(0)
+(1-\lambda)\,\max\left\{\sum_{i\in\I}\F_i(0)-1,\;0\right\},
\quad\lambda\in[0,1].
\end{equation}
Then there exists $\lambda^*\in[0,1]$ such that the coefficients
$\{p_\I(\lambda^*)\}_{\I\in\M}$ satisfy
\eqref{eq:JE_parameters_1} and~\eqref{eq:JE_parameters_2}.
\end{proposition}

\begin{proof}
Assume~\eqref{eq:JECondition} and define, for each $\I\in\M$,
\[
U_{\I}=\min_{i\in\I}\F_i(0),
\qquad
L_{\I}=\max\left\{\sum_{i\in\I}\F_i(0)-1,\;0\right\}.
\]
These are the Fr\'{e}chet upper and lower bounds for $p_\I$: $U_\I$ is the maximum
value consistent with~\eqref{eq:JE_parameters_1} individually, and $L_\I$ is the
minimum imposed by the bivariate Fr\'{e}chet lower bound applied
to~\eqref{eq:JEProbFaces}. The inequality $U_\I\ge L_\I$ holds because
$\min_{i\in\I}\F_i(0)\ge\sum_{i\in\I}\F_i(0)-1$ reduces to
$\max_{i\in\I}\F_i(0)\le1$, which is immediate.

Substituting the linear interpolation $p_\I(\lambda)=\lambda U_\I+(1-\lambda)L_\I$
into~\eqref{eq:JE_parameters_1} and~\eqref{eq:JE_parameters_2} and solving for
$\lambda$ yields the requirement
\begin{equation}
\label{eq:JE_trivariate_lambda}
\ell\;\le\;\lambda\;\le\;u_i\quad\text{for all }i\in\N,
\end{equation}
where
\[
\ell=\max\left\{
\frac{\sum_{i\in\N}\F_i(0)-1-\sum_{\I\in\M}L_{\I}}
     {\sum_{\I\in\M}(U_{\I}-L_{\I})},\;0\right\},
\qquad
u_i=\frac{\F_i(0)-\sum_{\I\in\M_i}L_{\I}}
         {\sum_{\I\in\M_i}(U_{\I}-L_{\I})}.
\]
We show that $\ell\le u_i\le1$ for every $i\in\N$, so that a valid $\lambda^*$
exists.

\smallskip
\noindent\textit{Non-negativity of $u_i$.}
For $n=3$ and each $i\in\N$, the two sets $\I\in\M_i$ each have $|\I|=2$, so
$L_\I=\max\{\F_i(0)+\F_k(0)-1,0\}$ for the other element $k\in\I\setminus\{i\}$.
The numerator $\F_i(0)-\sum_{\I\in\M_i}L_\I$ takes one of the following values
according to how many $L_\I$ are positive:
$\F_i(0)$, $1-\F_k(0)$ for $k\in\N\setminus\{i\}$, or
$2-\sum_{j\in\N}\F_j(0)$.
Each of these is non-negative: the first two trivially, and the last by
condition~\eqref{eq:JECondition}. Hence $u_i\ge0$.

\smallskip
\noindent\textit{Upper bound $u_i\le1$.}
The bound $u_i\le1$ is equivalent to $\F_i(0)\le\sum_{\I\in\M_i}U_\I$. For
$n=3$ and fixed $i$, $\sum_{\I\in\M_i}U_\I=\min(\F_i(0),\F_j(0))+\min(\F_i(0),\F_k(0))$
for $\{j,k\}=\N\setminus\{i\}$. This sum falls below $\F_i(0)$ only when
$\F_i(0)>\F_j(0)+\F_k(0)$, i.e.\ when $\F_i(0)$ strictly dominates both others.
Since $\F_j(0)>\F_i(0)$ and $\F_i(0)>\F_k(0)+\F_j(0)$ cannot hold simultaneously,
the condition $\sum_{\I\in\M_i}U_\I<\F_i(0)$ fails for at most one index $i$.
The minimum $\min_{i\in\N}u_i$ is therefore governed by the remaining indices, for
which $u_i\le1$, so $\min_{i\in\N}u_i\le1$.

\smallskip
\noindent\textit{Interval non-emptiness: $\ell\le\min_{i\in\N}u_i$.}
Since $\sum_{\I\in\M_i}(U_\I-L_\I)\le\sum_{\I\in\M}(U_\I-L_\I)$ and the
numerator of $u_i$ is non-negative,
\[
u_i\;\ge\;\frac{\F_i(0)-\sum_{\I\in\M_i}L_{\I}}{\sum_{\I\in\M}(U_{\I}-L_{\I})}.
\]
Using $\sum_{j\in\N\setminus\{i\}}\F_j(0)-1-L_{\N\setminus\{i\}}\le0$ (a
consequence of $L_{\N\setminus\{i\}}\ge\sum_{j\ne i}\F_j(0)-1$) and
$\sum_{\I\in\M}L_\I=\sum_{\I\in\M_i}L_\I+L_{\N\setminus\{i\}}$, we obtain
\begin{align*}
u_i
&\ge\frac{\F_i(0)-\sum_{\I\in\M_i}L_{\I}+\sum_{j\in\N\setminus\{i\}}\F_j(0)
         -1-L_{\N\setminus\{i\}}}{\sum_{\I\in\M}(U_{\I}-L_{\I})}
 =\frac{\sum_{j\in\N}\F_j(0)-1-\sum_{\I\in\M}L_{\I}}
       {\sum_{\I\in\M}(U_{\I}-L_{\I})}.
\end{align*}
Since this lower bound is precisely $\max\{\cdot,0\}$-truncated to give $\ell$,
we have $u_i\ge\ell$ for every $i\in\N$. Therefore, $\lambda^*$ exists in
$[\ell,\,\min_{i\in\N}u_i]\subseteq[0,1]$.
\end{proof}

%+++++++++++++++++++++++++++++++++++++++++++++++++++
\begin{example}
\label{ex:JE_canonical_example}
Let $n=3$, so that $\M=\bigl\{\{1,2\},\{1,3\},\{2,3\}\bigr\}$. Suppose
$\F_i(x_i)=\tfrac{1}{2}(1-x_i)$ for $x_i\in[0,1]$ and $i\in\N$, corresponding
to $X_i\sim\mathrm{Uniform}[0,1]$ marginals. Then $\sum_{i\in\N}\F_i(0)=3/2$,
so condition~\eqref{eq:JECondition} is satisfied. Applying
Proposition~\ref{prop:JE_trivariate} with $U_\I=\min_{i\in\I}\F_i(0)=1/2$ and
$L_\I=\max\{\sum_{i\in\I}\F_i(0)-1,0\}=0$ for every $\I\in\M$ gives
\[
p_{\I}(\lambda)=\frac{\lambda}{2},
\qquad
\lambda\in\left[\frac{1}{3},\,\frac{1}{2}\right].
\]
At $\lambda=1/3$, equality holds in~\eqref{eq:JE_parameters_2}, so the mass at
the origin vanishes (see Remark~\ref{rm:JE_Canonical_Remark}-(4)). At
$\lambda=1/2$, equality holds in~\eqref{eq:JE_parameters_1} for every $i$, so
all axes carry zero mass (see Remark~\ref{rm:JE_Canonical_Remark}-(3)).
Intermediate values of $\lambda$ distribute mass across both the faces and the
axes.

Now choose the copulas
\[
C_{12}(u_1,u_2)=u_1u_2,\quad
C_{13}(u_1,u_3)=\min\{u_1,u_3\},\quad
C_{23}(u_2,u_3)=\max\{u_2+u_3-1,0\},
\]
so that the three faces exhibit, respectively, independence, co-monotonicity, and
counter-monotonicity. By Corollary~\ref{cor:JE_Canonical_CDF_phi} and the
identity $\F_i(x_i)/\F_i(0)=1-x_i$ for $x_i\in[0,1]$, the joint CDF is
\[
F_{\N}(x_{\N})=\begin{cases}
{\displaystyle
1-\frac{1}{2}\sum_{i\in\N}(1-x_i)
+\frac{\lambda}{2}(1-x_1)(1-x_2)}\\[6pt]
{\displaystyle\quad
+\,\frac{\lambda}{2}\min\{1-x_1,1-x_3\}
+\frac{\lambda}{2}\max\{1-x_2-x_3,0\}},
& x_{\N}\in[0,1]^3,\\[8pt]
0, & \text{otherwise.}
\end{cases}
\]
The support of $X_{\N}$ is illustrated in Figure~\ref{fig:JE_canonical_ex}.

\begin{figure}[htbp]
\centering
\includegraphics[width=0.75\textwidth]{JE_canonical.pdf}
\caption{Support of the canonical JE random vector $X_{\N}$ with copulas
$C_{12}$, $C_{13}$, and $C_{23}$.}
\label{fig:JE_canonical_ex}
\end{figure}

The Pearson correlations are computed by observing that, on face $E_{\I}$,
$X_i|\text{face}\sim\mathrm{Uniform}[0,1]$, and contributions from all faces
not containing both components vanish. This gives $\E[X_i]=1/4$,
$\mathrm{Var}(X_i)=5/48$, and cross-moments
$\E[X_1X_2]=\lambda/8$ (independence),
$\E[X_1X_3]=\lambda/6$ (co-monotone, $X_1=X_3$ a.s.\ on face $E_{\{1,3\}}$),
$\E[X_2X_3]=\lambda/12$ (counter-monotone, $X_2+X_3=1$ a.s.\ on face $E_{\{2,3\}}$),
yielding
\[
\rho_{12}=\frac{6\lambda-3}{5},\qquad
\rho_{13}=\frac{8\lambda-3}{5},\qquad
\rho_{23}=\frac{4\lambda-3}{5}.
\]
For example,
\begin{align*}
\lambda=\tfrac{1}{3}:\quad
&\rho_{12}=-\tfrac{1}{5},\quad
\rho_{13}=-\tfrac{1}{15},\quad
\rho_{23}=-\tfrac{1}{3},\\
\lambda=\tfrac{5}{12}:\quad
&\rho_{12}=-\tfrac{1}{10},\quad
\rho_{13}=\tfrac{1}{15},\quad
\rho_{23}=-\tfrac{4}{15},\\
\lambda=\tfrac{1}{2}:\quad
&\rho_{12}=0,\quad
\rho_{13}=\tfrac{1}{5},\quad
\rho_{23}=-\tfrac{1}{5}.
\end{align*}
All three correlations increase linearly in $\lambda$, reflecting the diminishing
influence of axis mass (which induces strong negative dependence through the mutual
exclusion effect) as face mass grows. At $\lambda=1/2$, all mass is on the faces
and the correlations are determined entirely by the copulas $C_\I$.

\medskip
This example highlights the flexibility of the canonical JE framework. Although
the global structure enforces a form of counter-monotonicity through the JE
constraint, the model accommodates a wide range of local dependence structures
through the joint specification of the copulas $C_{\I}$ and coefficients $p_{\I}$.
\end{example}

Although the canonical JE framework covers a wide range of modelling choices,
the dependence structure within each face $E_{\I}$ is constrained by the marginals
through the ratios $\F_i(x_i)/\F_i(0)$. The copulas $C_\I$ may be chosen freely,
but the distributional mass on each face is spread over the range of the original
marginals. To afford greater control over the effective range of the face
distribution, we introduce a generalization that replaces the ratio
$\F_i(x_i)/\F_i(0)$ by a prescribed transformation $G_i(\F_i(x_i))$. We call
the resulting class the \emph{generalised JE} or G-JE.

\begin{definition}
\label{def:g-JE}
For each $i\in\N$, let $G_i:[0,1]\to[0,1]$ be a non-decreasing, left-continuous
function satisfying $G_i(u)=0$ for $u\in[0,a_i]$ and $G_i(u)=1$ for
$u\in[b_i,1]$, where $0\le a_i<b_i\le\F_i(0)$. The boundary condition
$b_i\le\F_i(0)$ ensures $G_i(\F_i(0))=1$, which guarantees that the face
probability at $x_\I=0$ equals $p_\I$ and that the origin mass formula is
unchanged. The composition $G_i\circ\F_i$ is right-continuous in $x_i$ since
$\F_i$ is right-continuous and non-increasing, and left-continuity of $G_i$
ensures $G_i(\F_i(x_n))\to G_i(\F_i(x))$ as $x_n\downarrow x$.

A random vector $X_{\N}$ is said to be \emph{G-JE} if, for each $\I\in\M$,
\begin{equation}
\label{eq:G-JEProbFaces}
\P\left(X_i>x_i,\;i\in \I;\;X_j=0,\;j\in\N\setminus \I\right)
= p_{\I}\,C_{\I}\left(G_i\left(\F_i(x_i)\right),\;i\in\I\right),
\quad x_{\I}\in[0,\infty)^{|\I|};
\end{equation}
for each $i\in\N$,
\begin{equation}
\label{eq:G-JEProbAxes}
\P\left(X_i>x_i;\;X_j=0,\;j\in\N\setminus\{i\}\right)
=\F_i(x_i)-\left(\sum_{\I\in \M_i}p_{\I}\right)G_i\left(\F_i(x_i)\right),
\quad x_i\in[0,\infty);
\end{equation}
and the origin probability is as in~\eqref{eq:JEProbOrigin}. The coefficients
$p_\I\ge0$ satisfy condition~\eqref{eq:JE_parameters_2} together with
\begin{equation}
\label{eq:G-JE_parameters_1}
\sum_{\I\in \M_i}p_{\I}\le G_i^*, \qquad i\in\N,
\end{equation}
where
\begin{equation}
\label{eq:G-JE_G_star}
G_i^*=\inf_{\substack{0\le x_i<y_i\\
G_i(\F_i(x_i))>G_i(\F_i(y_i))}}
\frac{\F_i(x_i)-\F_i(y_i)}{G_i\left(\F_i(x_i)\right)-G_i\left(\F_i(y_i)\right)}.
\end{equation}
\end{definition}

The quantity $G_i^*$ is the smallest local ratio of the decrease in $\F_i$ to
the corresponding decrease in $G_i\circ\F_i$; it controls how much face mass
the $i$-th axis can support while remaining a valid non-negative, non-increasing
function of $x_i$. The canonical JE is recovered as the special case
$G_i(u)=u/\F_i(0)$, for which $G_i^*=\F_i(0)$ and condition~\eqref{eq:G-JE_parameters_1}
reduces to~\eqref{eq:JE_parameters_1}.

\begin{theorem}
\label{thm:G_JE_characterization}
Let $F_1,\dots,F_n$ be marginal distributions on $[0,\infty)$ and let
$G_1,\dots,G_n$ be functions as in Definition~\ref{def:g-JE}. A G-JE random
vector $X_{\N}$ with marginals $F_1,\dots,F_n$ exists if and only if
\begin{equation}
\label{eq:G-JECondition}
\sum_{i\in\N}\F_i(0)-1
\le \min\left\{\frac{n-2}{n-1}\sum_{i\in\N}G_i^*,\;
\sum_{i\in\N}G_i^*-\max_{i\in\N}G_i^*\right\}.
\end{equation}
\end{theorem}

\begin{proof}
\textit{Validity of the construction.}
Since each $G_i$ is left-continuous with $G_i(u)=0$ for $u\le a_i$ and
$G_i(\F_i(0))=1$, the face probabilities in~\eqref{eq:G-JEProbFaces} define
valid DDFs on $E_{\I}$: at $x_\I=0$ they give $p_\I C_\I(1,\dots,1)=p_\I$, and
they are non-increasing in each $x_i$ and right-continuous.

\textit{Axis probabilities.}
Define $h_i(u)=u-s_i G_i(u)$ for $u\in[0,\F_i(0)]$, where
$s_i=\sum_{\I\in\M_i}p_\I$. The axis probability for component $i$ equals
$h_i(\F_i(x_i))$. For this to define a valid (non-negative, non-increasing in
$x_i$) survival function on the $i$-th axis, $h_i$ must be non-decreasing as a
function of $u$ on $[0,\F_i(0)]$. Since $G_i(0)=0$, we have $h_i(0)=0$;
non-decreasingness of $h_i$ then ensures non-negativity. For any
$u>v$ with $G_i(u)>G_i(v)$, non-decreasingness requires
$u-v\ge s_i(G_i(u)-G_i(v))$, i.e.\ $s_i\le(u-v)/(G_i(u)-G_i(v))$. Taking the
infimum over all such pairs yields condition~\eqref{eq:G-JE_parameters_1}, with
$G_i^*$ as defined in~\eqref{eq:G-JE_G_star}. Moreover, since $G_i(\F_i(0))=1$,
taking $y_i\to\infty$ in~\eqref{eq:G-JE_G_star} shows $G_i^*\le\F_i(0)$,
ensuring non-negativity also at $x_i=0$.

\textit{Origin probability.}
Since $G_i(\F_i(0))=1$ for each $i$, the total mass on the axes equals
$\sum_{i\in\N}[\F_i(0)-s_i]$ and the total face mass is $\sum_{\I\in\M}p_\I$.
Subtracting from unity gives the origin mass $1-\sum_{i\in\N}\F_i(0)+\sum_{\I\in\M}(|\I|-1)p_\I$,
which coincides with~\eqref{eq:JEProbOrigin} and is non-negative if and only
if~\eqref{eq:JE_parameters_2} holds.

\textit{Necessity and sufficiency of~\eqref{eq:G-JECondition}.}
The construction yields a valid G-JE distribution if and only if there exist
$p_\I\ge0$ simultaneously satisfying~\eqref{eq:G-JE_parameters_1}
and~\eqref{eq:JE_parameters_2}. Applying the same linear programming argument as
in~\eqref{sys:original}--\eqref{sys:dual}, with $G_i^*$ in place of $\F_i(0)$,
the primal maximum of $\sum_{\I\in\M}(|\I|-1)p_\I$ subject to
$p_\I\ge0$ and~\eqref{eq:G-JE_parameters_1} equals, by strong duality,
\[
\min\left\{\frac{n-2}{n-1}\sum_{i\in\N}G_i^*,\;
\sum_{i\in\N}G_i^*-\max_{i\in\N}G_i^*\right\}.
\]
Feasible $p_\I$ satisfying~\eqref{eq:JE_parameters_2} exist if and only if this
maximum is at least $\sum_{i\in\N}\F_i(0)-1$, which is precisely
condition~\eqref{eq:G-JECondition}.
\end{proof}


%+++++++++++++++++++++++++++++++++++++++++++++++++++
\begin{remark}
\label{rm:G-JE_Remark}
\mbox{}
\begin{itemize}

\item[(1)]
The G-JE is strictly more flexible than canonical JE: the parameters $a_i$ and
$b_i$ restrict the effective support on each face $E_\I$ to a sub-range of the
marginal support, while the functions $G_i$ introduce distributional distortions
within that range, all without disrupting the global JE constraint. The cost is a
more stringent existence condition~\eqref{eq:G-JECondition}, owing to the
inequality $G_i^*\le\F_i(0)$ established in the proof of
Theorem~\ref{thm:G_JE_characterization}. When $\F_i$ attains both boundary
values $a_i$ and $b_i$, this bound sharpens further: if there exist
$x_i^{(1)}<x_i^{(2)}$ such that $\F_i(x_i^{(1)})=b_i$ and
$\F_i(x_i^{(2)})=a_i$, then
\[
G_i^*\le\frac{b_i-a_i}{G_i(b_i)-G_i(a_i)}=b_i-a_i\le\F_i(0),
\]
so the admissible range for $\{p_\I\}_{\I\in\M}$ is correspondingly smaller
than in the canonical case.

\item[(2)]
Suppose $\F_i$ is continuous at some $x_i^*\in(0,\infty)$ while
$G_i\circ\F_i$ has a left-jump there, i.e.\
$c:=\lim_{x\nearrow x_i^*}G_i(\F_i(x))-G_i(\F_i(x_i^*))>0$.
Since $G_i$ is left-continuous and $\F_i$ is right-continuous, such a
left-jump in $G_i\circ\F_i$ arises precisely when $G_i$ has a
right-discontinuity at $\F_i(x_i^*)$: as $x_i^{(k)}\nearrow x_i^*$, one has
$\F_i(x_i^{(k)})\searrow\F_i(x_i^*)$ from above, so $G_i(\F_i(x_i^{(k)}))$
converges to the right-limit $G_i(\F_i(x_i^*)^+)=G_i(\F_i(x_i^*))+c$.
Taking $y_i=x_i^*$ and $x_i=x_i^{(k)}\nearrow x_i^*$
in~\eqref{eq:G-JE_G_star}: continuity of $\F_i$ gives
$\F_i(x_i^{(k)})-\F_i(x_i^*)\to0$, while the denominator converges to $c>0$.
Hence $G_i^*=0$, and condition~\eqref{eq:G-JE_parameters_1} forces
\[
p_{\I}=0\quad\text{for all }\I\in\M_i.
\]
This constraint is necessary: if $s_i=\sum_{\I\in\M_i}p_\I>0$, then the
axis survival function $h_i(x_i)=\F_i(x_i)-s_i G_i(\F_i(x_i))$ satisfies
$\lim_{x\nearrow x_i^*}h_i(x)=h_i(x_i^*)-s_i c<h_i(x_i^*)$,
so the left-limit of $h_i$ at $x_i^*$ is strictly less than its value,
contradicting the required non-increasing property. Consequently, wherever
$\F_i$ is continuous, $G_i\circ\F_i$ must be continuous as well.

\item[(3)]
If $G_i\circ\F_i$ is constant on some sub-interval while $\F_i$ is strictly
decreasing there, the denominator in~\eqref{eq:G-JE_G_star} vanishes for the
corresponding pairs $(x_i,y_i)$. These pairs are excluded from the infimum by
the constraint $G_i(\F_i(x_i))>G_i(\F_i(y_i))$, so no difficulty arises.

\item[(4)]
The CDF of a G-JE random vector $X_\N$ takes the same form as the canonical
expression~\eqref{eq:JE_Canonical_CDF}, with each copula argument
$\F_i(x_i)/\F_i(0)$ replaced by $G_i(\F_i(x_i))$.

\item[(5)]
Remark~\ref{rm:JE_Canonical_Remark}-(1) carries over to G-JE without
modification. Remark~\ref{rm:JE_Canonical_Remark}-(2) extends with the
additional requirement $G_i^*=\F_i(0)=1$ for all $i\in\I$, which forces
those components to reduce to the canonical case.
Remark~\ref{rm:JE_Canonical_Remark}-(4) carries over without modification.

Analogously to Remark~\ref{rm:JE_Canonical_Remark}-(3), the $i$-th axis
carries zero mass if and only if equality holds in~\eqref{eq:G-JE_parameters_1}:
$\sum_{\I\in\M_i}p_\I=G_i^*$. Substituting
into~\eqref{eq:G-JEProbAxes} gives $G_i(\F_i(x_i))=\F_i(x_i)/G_i^*$ for
all $x_i\ge0$; evaluating at $x_i=0$ and using $G_i(\F_i(0))=1$ yields
$G_i^*=\F_i(0)$ and $G_i(u)=u/\F_i(0)$, reducing component $i$ to the
canonical case. Consequently, for $n=3$, requiring zero mass on all three
axes forces $G_i^*=\F_i(0)$ for every $i\in\N$, and the system
$\sum_{\I\in\M_i}p_\I=\F_i(0)$ reduces to its canonical counterpart with
unique solution
\[
p_\I=\frac{\sum_{i\in\I}\F_i(0)-\F_j(0)}{2},\qquad j\in\N\setminus\I,
\]
valid when $\{\F_i(0)\}_{i\in\N}$ satisfies the triangle inequality
(see Remark~\ref{rm:JE_Canonical_Remark}-(3)).

The analog of Remark~\ref{rm:JE_Canonical_Remark}-(5), however, does not
extend to G-JE. In the canonical setting, the tight condition
$\sum_{i\in\N}\F_i(0)=n-1$ implies $(n-1)a_0+(n-2)\sum_{i\in\N}a_i\le0$,
forcing all axis and origin masses to zero. For G-JE, the existence
condition~\eqref{eq:G-JECondition} becomes tight when
$\sum_{i\in\N}\F_i(0)-1$ equals the right-hand side of~\eqref{eq:G-JECondition},
a quantity governed by $\{G_i^*\}$ rather than $\{\F_i(0)\}$. Since each axis
carries a minimum mass of $\F_i(0)-s_i\ge\F_i(0)-G_i^*\ge0$ regardless of
the choice of $\{p_\I\}$, a tight G-JE condition does not in general force
the axis masses to vanish. For instance, with $n=3$, $\F_i(0)=\tfrac{1}{2}$
and $G_i^*=\tfrac{1}{3}$ for all $i\in\N$,
condition~\eqref{eq:G-JECondition} holds with equality, yet each axis
carries a minimum mass of $\F_i(0)-G_i^*=\tfrac{1}{6}>0$.
\end{itemize}
\end{remark}

Analogously to Proposition~\ref{prop:JE_trivariate}, the admissible coefficients
$\{p_\I\}_{\I\in\M}$ for a trivariate G-JE random vector admit a one-parameter
linear interpolation. The role of $\F_i(0)$ in the upper bounds $U_\I$ is now
played by $G_i^*$, while the lower bounds $L_\I$ are adjusted to reflect
constraint~\eqref{eq:G-JE_parameters_1} in place of~\eqref{eq:JE_parameters_1}.

\begin{proposition}
\label{prop:G-JE_trivariate}
Let $n=3$ and suppose that condition~\eqref{eq:G-JECondition} holds. For each
$\I\in\M$ with complementary index $j\in\N\setminus\I$, define
\begin{equation}
\label{eq:G-JE_trivariate_P}
p_{\I}(\lambda)
=\lambda\,\min_{i\in\I}G_i^*
+(1-\lambda)\,\max\left\{\sum_{k\in\N}\F_k(0)-G_j^*-1,\;0\right\},
\quad\lambda\in[0,1].
\end{equation}
Then there exists $\lambda^*\in[0,1]$ such that $\{p_\I(\lambda^*)\}_{\I\in\M}$
satisfies both~\eqref{eq:G-JE_parameters_1} and~\eqref{eq:JE_parameters_2}.
\end{proposition}

\begin{proof}
Assume~\eqref{eq:G-JECondition} holds and define, for each $\I\in\M$ with
$j\in\N\setminus\I$,
\[
U_\I=\min_{i\in\I}G_i^*,\qquad
L_\I=\max\left\{\sum_{k\in\N}\F_k(0)-G_j^*-1,\;0\right\}.
\]
Here $U_\I$ is the maximum value of $p_\I$ consistent
with~\eqref{eq:G-JE_parameters_1} when all other face masses vanish, and
$L_\I$ is the lower bound on $p_\I$ derived from the combined
constraints~\eqref{eq:G-JE_parameters_1} and~\eqref{eq:JE_parameters_2}: for
$n=3$ and $\I=\{i,k\}$ with $j\in\N\setminus\I$, constraint~\eqref{eq:G-JE_parameters_1}
for index $j$ gives $p_{\{i,j\}}+p_{\{k,j\}}\le G_j^*$; substituting
into~\eqref{eq:JE_parameters_2} yields $p_\I\ge\sum_{l\in\N}\F_l(0)-G_j^*-1$,
hence $L_\I$ as stated.

The inequality $U_\I\ge L_\I$ holds by condition~\eqref{eq:G-JECondition}: for
any $j\in\N\setminus\I$,
\[
\min_{i\in\I}G_i^*+G_j^*
\ge\sum_{l\in\N}G_l^*-\max_{l\in\N}G_l^*
\ge\sum_{l\in\N}\F_l(0)-1,
\]
where the first inequality follows because $\max_l G_l^*\ge G_l^*$ for every $l$,
and the second from~\eqref{eq:G-JECondition}; rearranging gives
$U_\I\ge\sum_{l\in\N}\F_l(0)-G_j^*-1$.

Substituting $p_\I(\lambda)=\lambda U_\I+(1-\lambda)L_\I$
into~\eqref{eq:G-JE_parameters_1} and~\eqref{eq:JE_parameters_2} and solving
for $\lambda$ yields the requirement
\begin{equation}
\label{eq:G-JE_trivariate_lambda}
\ell\;\le\;\lambda\;\le\;u_i\quad\text{for all }i\in\N,
\end{equation}
where
\[
\ell=\max\left\{
\frac{\sum_{k\in\N}\F_k(0)-1-\sum_{\I\in\M}L_{\I}}
     {\sum_{\I\in\M}(U_{\I}-L_{\I})},\;0\right\},
\qquad
u_i=\frac{G_i^*-\sum_{\I\in\M_i}L_{\I}}
         {\sum_{\I\in\M_i}(U_{\I}-L_{\I})}.
\]
We show that $0\le\ell\le u_i\le1$ for every $i\in\N$.

\smallskip
\noindent\textit{Non-negativity of $u_i$.}
For fixed $i\in\N$, the numerator $G_i^*-\sum_{\I\in\M_i}L_\I$ takes one of
the following values according to how many terms in $\sum_{\I\in\M_i}L_\I$ are
strictly positive: $G_i^*$ (both zero); $G_i^*+G_j^*-(\sum_k\F_k(0)-1)$ for
$j\in\N\setminus\{i\}$ (exactly one positive); or
$\sum_k G_k^*-2(\sum_k\F_k(0)-1)$ (both positive). All three are non-negative
by condition~\eqref{eq:G-JECondition}: the first trivially; the second since
$\sum_k\F_k(0)-1\le\sum_k G_k^*-\max_k G_k^*\le G_i^*+G_j^*$ for any
$j\ne i$; the third since $\sum_k\F_k(0)-1\le\tfrac{1}{2}\sum_k G_k^*$ (first
term of~\eqref{eq:G-JECondition} with $n=3$). Hence $u_i\ge0$.

\smallskip
\noindent\textit{Upper bound $u_i\le1$.}
The bound $u_i\le1$ is equivalent to
$G_i^*\le\sum_{\I\in\M_i}U_\I=\min(G_i^*,G_j^*)+\min(G_i^*,G_k^*)$ for
$\{j,k\}=\N\setminus\{i\}$. This sum falls below $G_i^*$ only if
$G_i^*>G_j^*+G_k^*$, a strict dominance condition that can hold for at most
one index $i$ simultaneously. Hence $\min_{i\in\N}u_i\le1$.

\smallskip
\noindent\textit{Non-emptiness: $\ell\le\min_{i\in\N}u_i$.}
Since $\sum_{\I\in\M_i}(U_\I-L_\I)\le\sum_{\I\in\M}(U_\I-L_\I)$ and the
numerator of $u_i$ is non-negative, $u_i\ge(G_i^*-\sum_{\I\in\M_i}L_\I)/
\sum_{\I\in\M}(U_\I-L_\I)$.
By definition, $L_{\N\setminus\{i\}}=\max\{\sum_k\F_k(0)-G_i^*-1,0\}\ge
\sum_k\F_k(0)-G_i^*-1$, so
$\sum_k\F_k(0)-G_i^*-1-L_{\N\setminus\{i\}}\le0$. Writing
$\sum_{\I\in\M}L_\I=\sum_{\I\in\M_i}L_\I+L_{\N\setminus\{i\}}$ and adding
the non-positive term $\sum_k\F_k(0)-G_i^*-1-L_{\N\setminus\{i\}}$ to the
numerator:
\begin{align*}
u_i
&\ge\frac{G_i^*-\sum_{\I\in\M_i}L_\I
         +\sum_k\F_k(0)-G_i^*-1-L_{\N\setminus\{i\}}}
        {\sum_{\I\in\M}(U_\I-L_\I)}
=\frac{\sum_k\F_k(0)-1-\sum_{\I\in\M}L_\I}{\sum_{\I\in\M}(U_\I-L_\I)}.
\end{align*}
Since this lower bound equals the untruncated part of $\ell$, we have
$u_i\ge\ell$ for every $i\in\N$. Therefore $\lambda^*$ exists in
$[\ell,\,\min_{i\in\N}u_i]\subseteq[0,1]$.
\end{proof}

%+++++++++++++++++++++++++++++++++++++++++++++++++++
In the next example, we illustrate several basic choices of the functions $G_i$
and the rich G-JE structures they generate.

\begin{example}
\label{ex:G-JE_example}
\mbox{}
\begin{itemize}

\item[(1)]
Consider the linear functions
\[
G_i(u)=\frac{u-a_i}{b_i-a_i},\quad u\in[a_i,b_i],
\]
extended by $G_i(u)=0$ for $u\in[0,a_i]$ and $G_i(u)=1$ for $u\in[b_i,1]$.
The boundary conditions $G_i(a_i)=0$ and $G_i(b_i)=1$ are satisfied by
construction, and uniquely determine the slope $1/(b_i-a_i)$. Since $G_i$ is
linear on $[a_i,b_i]$, the ratio $(\F_i(x_i)-\F_i(y_i))/
(G_i(\F_i(x_i))-G_i(\F_i(y_i)))$ equals $b_i-a_i$ whenever both
$\F_i(x_i),\F_i(y_i)\in[a_i,b_i]$, and is larger otherwise.
Hence $G_i^*=b_i-a_i$, provided $\F_i$ attains at least one value in
$(a_i,b_i)$. Setting $a_i=0$ and $b_i=\F_i(0)$ gives $G_i(u)=u/\F_i(0)$ and
$G_i^*=\F_i(0)$, recovering the canonical JE.

Now take $a_i>0$ and $b_i=\F_i(0)$, and assume each $\F_i$ attains at least one
value in $[a_i,\F_i(0))$. Then $G_i^*=\F_i(0)-a_i$, and
condition~\eqref{eq:G-JECondition} becomes
\[
\sum_{i\in\N}\F_i(0)-1
\le\min\left\{
\frac{n-2}{n-1}\sum_{i\in\N}\bigl(\F_i(0)-a_i\bigr),\;
\sum_{i\in\N}\bigl(\F_i(0)-a_i\bigr)-\max_{i\in\N}\bigl(\F_i(0)-a_i\bigr)
\right\}.
\]
Writing $\F_i(0)-a_i = G_i^*$ and expanding each case:

\noindent\textit{Case 1} (symmetric regime, first term binding): the condition
reduces to
\[
\sum_{i\in\N}\F_i(0)+(n-2)\sum_{i\in\N}a_i\le n-1.
\]
Compared with~\eqref{eq:JECondition}, the extra term $(n-2)\sum_{i\in\N}a_i$
reflects the additional restriction imposed by truncating the face distributions
at $a_i$ from below.

\noindent\textit{Case 2} (asymmetric regime, second term binding): the condition
becomes
\[
\sum_{i\in\N}a_i+\max_{i\in\N}\bigl(\F_i(0)-a_i\bigr)\le 1,
\]
a constraint that primarily controls the balance between the lower bounds
$\{a_i\}$ and the largest effective face mass.

Setting $a_i=1-q_i$ for parameters $q_i\in[F_i(0),1]$ (so that
$G_i^*=q_i-F_i(0)\ge0$) and using $\F_i(0)=1-F_i(0)$, these two conditions
may equivalently be written as
\begin{align*}
\sum_{i\in\N}(1-q_i)
+\frac{\sum_{i\in\N}(q_i-F_i(0))}{n-1}&\le 1,\\
\sum_{i\in\N}(1-q_i)+\max_{i\in\N}(q_i-F_i(0))&\le 1.
\end{align*}
The second inequality is precisely the necessary and sufficient condition for the
existence of the Tail-ME random vector with probability vector $(q_1,\dots,q_n)$
studied in~\cite{Cheung2017}. The first inequality additionally ensures that the
G-JE variant concentrates its face mass on the lower portion of the box
$B=\{x_{\N}\in[0,\infty)^n:0\le x_i\le F_i^{-1}(q_i)\}$, rather than over the
full $n$-cube $B$.

As a concrete illustration, modify Example~\ref{ex:JE_canonical_example} by
setting $a_i=1/8$ (and retaining $b_i=\F_i(0)=1/2$) for all $i\in\N$.
Then $G_i^*=3/8$, and one verifies that~\eqref{eq:G-JECondition} holds. The
transformed argument is
\[
G_i\left(\F_i(x_i)\right)
=\frac{(1-x_i)/2-1/8}{3/8}
=\left(1-\tfrac{4}{3}x_i\right)_+,
\quad x_i\in[0,1],
\]
where $(t)_+=\max\{t,0\}$ denotes the positive part. Applying
Proposition~\ref{prop:G-JE_trivariate} with $U_\I=3/8$ and $L_\I=1/8$ for
every $\I\in\M$ gives
\[
p_\I(\lambda)
=\tfrac{3}{8}\lambda+\tfrac{1}{8}(1-\lambda)
=\tfrac{\lambda}{4}+\tfrac{1}{8},
\quad\lambda\in\left[\tfrac{1}{6},\,\tfrac{1}{4}\right].
\]
Retaining the copulas $C_{12}$, $C_{13}$, $C_{23}$ from
Example~\ref{ex:JE_canonical_example} and substituting the transformed arguments,
the CDF of $X_\N$ is
\begin{align*}
F_{\N}(x_{\N})=\begin{cases}
{\displaystyle 1-\frac{1}{2}\sum_{i\in\N}(1-x_i)}\\[6pt]
{\displaystyle
+\left(\frac{\lambda}{4}+\frac{1}{8}\right)\Bigl[
\left(1-\tfrac{4}{3}x_1\right)_{+}
\left(1-\tfrac{4}{3}x_2\right)_{+}}\\[4pt]
{\displaystyle\qquad
+\min\left\{\left(1-\tfrac{4}{3}x_1\right)_{+},
             \left(1-\tfrac{4}{3}x_3\right)_{+}\right\}}\\[4pt]
{\displaystyle\qquad
+\max\left\{\left(1-\tfrac{4}{3}x_2\right)_{+}
             +\left(1-\tfrac{4}{3}x_3\right)_{+}-1,\;0\right\}
\Bigr],}
& x_{\N}\in[0,1]^3,\\[6pt]
0, & \text{otherwise.}
\end{cases}
\end{align*}
The effective face support is confined to $[0,3/4)$ in each active coordinate,
strictly smaller than the $[0,1]$ range of the canonical JE; the support of
$X_\N$ is depicted in Figure~\ref{fig:G-JE_ex}.

\begin{figure}[htbp]
\centering
\includegraphics[width=0.75\textwidth]{G-JE.pdf}
\caption{Support of the G-JE random vector $X_\N$ with linear truncation
functions $G_i$ at level $a_i=1/8$, restricting the effective face support to
$[0,3/4)$ in each active coordinate.}
\label{fig:G-JE_ex}
\end{figure}

\item[(2)]
Consider the power functions
\[
G_i(u)=\left(\frac{u-a_i}{b_i-a_i}\right)^{\gamma},\quad u\in[a_i,b_i],
\]
with exponent $\gamma>0$, extended by $G_i(u)=0$ for $u\in[0,a_i]$ and
$G_i(u)=1$ for $u\in[b_i,1]$. This recovers the linear case at $\gamma=1$.
For general $\gamma$, the ratio in~\eqref{eq:G-JE_G_star} equals
$(b_i-a_i)^\gamma(u-v)/(u^\gamma-v^\gamma)$ (with $u = s+a_i$, $v = t+a_i$, $s>t\ge0$).
By the mean value theorem, $(u^\gamma-v^\gamma)/(u-v)=\gamma\xi^{\gamma-1}$ for
some $\xi\in(v,u)$, so the ratio equals $(b_i-a_i)^\gamma/(\gamma\xi^{\gamma-1})$.

\noindent\textit{Case $\gamma\in(0,1)$}: since $\gamma-1<0$, the ratio vanishes
as $\xi\to0^+$, which occurs when $v\to a_i^+$. Hence $G_i^*=0$ whenever $\F_i$
attains values in $(a_i,a_i+\varepsilon)$ for some $\varepsilon>0$.
Condition~\eqref{eq:G-JECondition} then reduces to $\sum_{i\in\N}\F_i(0)\le 1$,
the ME condition~\eqref{eq:MEcondition}: the strong concavity of $G_i$ near
$a_i$ forces all face probabilities to vanish, recovering ME as the only feasible
structure.

\noindent\textit{Case $\gamma>1$}: the ratio is minimised as $\xi\to b_i$
(since $\gamma-1>0$ and $\xi^{\gamma-1}$ is increasing), which occurs as
$v\to u^-$ with $u\to b_i$. Setting $a_i=0$ and $b_i=\F_i(0)$ and assuming
$\F_i$ attains values in $(b_i-\varepsilon,b_i)$ for some $\varepsilon>0$, one
obtains $G_i^*=\F_i(0)/\gamma$. Condition~\eqref{eq:G-JECondition} then becomes
\[
\sum_{i\in\N}\F_i(0)
\le\min\left\{
\frac{\gamma(n-1)}{\gamma(n-1)-(n-2)},\;
\frac{\gamma-\max_{i\in\N}\F_i(0)}{\gamma-1}
\right\}.
\]
Both terms are decreasing in $\gamma$ from the canonical JE bound of $n-1$ (at
$\gamma=1^+$) toward $1$ as $\gamma\to\infty$, confirming that larger exponents
impose progressively stronger constraints relative to the canonical case.
\end{itemize}
\end{example}

Another extension to JE comes as an $m$-exclusivity hierarchy in which JE corresponds to the
least restrictive level. Strengthening the exclusion condition of JE to smaller sub-groups of
components yields a nested family of conditions interpolating between JE and ME.

\begin{definition}
\label{def:m_exclusive}
Let $2\le m\le n$. A non-negative random vector $X_{\N}$ is said to be
\emph{$m$-exclusive} if
\[
\F_{I}(0_{I})=0,\text{ or equivalently, }
    \prod_{i\in\I}X_i=0\text{ almost surely},\quad\text{for all }I\subseteq\N\text{ with }|I|=m,
\]
that is, no $m$ components are simultaneously strictly positive almost surely.
The special cases $m=n$ and $m=2$ recover JE and ME, respectively.
\end{definition}

Since $\P(X_i>0,\,i\in J)\le \P(X_i>0,\,i\in I)$ for any $I\subseteq J$,
$m$-exclusivity immediately implies $d$-exclusivity for every $d\ge m$: any
set of $d>m$ components contains a sub-group of size $m$ with zero joint
positivity probability, hence itself has zero joint positivity probability.
Consequently, the $m$-exclusivity conditions form a decreasing chain:
\[
\text{ME} = 2\text{-exclusive}
\;\Rightarrow\;3\text{-exclusive}
\;\Rightarrow\;\cdots
\;\Rightarrow\;n\text{-exclusive} = \text{JE}.
\]

\begin{theorem}
\label{thm:m-exclusivity}
Let $F_1,\dots,F_n$ be marginal distributions on $[0,\infty)$ and let
$2\le m\le n$. An $m$-exclusive random vector $X_{\N}$ with these marginals
exists if and only if
\begin{equation}
\label{eq:m-exclusive_Condition}
\sum_{i\in\N}\F_i(0)\le m-1.
\end{equation}
\end{theorem}

\begin{proof}
\textit{Necessity.}
If $X_\N$ is $m$-exclusive, then no $m$ components are simultaneously positive,
so $\sum_{i\in\N}\mathds{1}_{\{X_i>0\}}\le m-1$ almost surely. Taking expectations
gives $\sum_{i\in\N}\F_i(0)\le m-1$.

\medskip
\noindent\textit{Sufficiency.}
Assume~\eqref{eq:m-exclusive_Condition} holds. The proof proceeds exactly as
that of Theorem~\ref{thm:JE_Characterisation_1}, with the face collection
$\M$ replaced by $\M^{(m)}=\{I\subset\N:2\le|I|\le m-1\}$. Restricting mass
to faces of size at most $m-1$ ensures that no $m$ components are ever
simultaneously positive, thereby enforcing $m$-exclusivity by construction.

The feasibility argument (Steps~1--4 of the proof of
Theorem~\ref{thm:JE_Characterisation_1}) applies to the linear program
with face collection $\M^{(m)}$. The dual constraint $\sum_{i\in\I}r_i\ge|\I|-1$
for $\I\in\M^{(m)}$ is satisfied at the symmetric vertex
$r^*=((m-2)/(m-1),\dots,(m-2)/(m-1))$, since $|\I|(m-2)/(m-1)\ge|\I|-1$ holds
for all $|\I|\le m-1$, and at each permutation of $(1,\dots,1,0)$. By strong
duality, the maximum of $\sum_{\I\in\M^{(m)}}(|\I|-1)p_\I$ equals
\[
\min\left\{\frac{m-2}{m-1}\sum_{i\in\N}\F_i(0),\;
\sum_{i\in\N}\F_i(0)-\max_{i\in\N}\F_i(0)\right\}.
\]
In both cases, this minimum is at least $\sum_{i\in\N}\F_i(0)-1$ under
condition~\eqref{eq:m-exclusive_Condition}: the second term yields the trivial
bound $\max_{i\in\N}\F_i(0)\le1$, while the first requires
$\sum_{i\in\N}\F_i(0)\le m-1$.
\end{proof}

\begin{corollary}
\label{cor:m_exclusive_G_star}
Let $F_1,\dots,F_n$ be marginal distributions on $[0,\infty)$, let $2\le m\le n$,
and let $G_1,\dots,G_n$ be functions as in Definition~\ref{def:g-JE}. A G
$m$-exclusive random vector $X_\N$ with marginals $F_1,\dots,F_n$ exists if and
only if
\begin{equation}
\label{eq:G-m-exclusive_Condition}
\sum_{i\in\N}\F_i(0)-1
\le\min\left\{
\frac{m-2}{m-1}\sum_{i\in\N}G_i^*,\;
\sum_{i\in\N}G_i^*-\max_{i\in\N}G_i^*
\right\}.
\end{equation}
\end{corollary}

\begin{proof}
The proof follows the argument of Theorem~\ref{thm:G_JE_characterization},
applied to the restricted face collection $\M^{(m)}=\{I\subset\N:2\le|I|\le m-1\}$
in place of $\M$, with $G_i^*$ in place of $\F_i(0)$ in the linear program, and
with $m$ taking the role of $n$ throughout the dual analysis.
\end{proof}
\begin{remark}
\label{rm:m_exclusive_canonical}
The $m$-exclusivity hierarchy admits a transparent interpretation within the
canonical JE framework.

\textit{Connection to face probabilities.}
A canonical JE random vector $X_\N$ is $m$-exclusive if and only if $p_\I=0$
for every $\I\in\M$ with $|\I|\ge m$. For the forward direction: if $X_\N$ is
$m$-exclusive and $\J\in\M$ with $|\J|\ge m$, choose any $I\subseteq\J$ with
$|I|=m$; then
\[
p_\J\le\sum_{\substack{\K\in\M\\\K\supseteq I}}p_\K
=\P(X_i>0,\;i\in I)=0.
\]
Conversely, if $p_\J=0$ for all $|\J|\ge m$, then for any $I\subseteq\N$ with
$|I|=m$:
\[
\P(X_i>0,\;i\in I)
=\sum_{\substack{\K\in\M\\\K\supseteq I,\,|\K|\ge m}}p_\K=0,
\]
confirming $m$-exclusivity. As $m$ decreases from $n$ to $2$, faces of
successively larger dimension are suppressed: at $m=n$ the full canonical JE
support is recovered; at $m=2$ all faces are zero ($p_\I=0$ for all $\I\in\M$),
consistent with Remark~\ref{rm:JE_Canonical_Remark}-(1), and mass concentrates
entirely on axes and origin, recovering ME.

\textit{Parallel structural properties.}
The results of Remark~\ref{rm:JE_Canonical_Remark} extend to $m$-exclusive
vectors with $\M$ replaced by $\M^{(m)}$ and~\eqref{eq:JECondition} replaced
by~\eqref{eq:m-exclusive_Condition}. Items~(1)--(4) carry over without
modification. For item~(5), suppose $\sum_{i\in\N}\F_i(0)=m-1$. Summing
\eqref{eq:JEProbAxes} over all $i\in\N$ at $x_i=0$ with $\M^{(m)}$ in place
of $\M$, and using $|\I|\le m-1$ for every $\I\in\M^{(m)}$, the same argument
as Remark~\ref{rm:JE_Canonical_Remark}-(5) yields
\[
(m-1)\,c_0+(m-2)\sum_{i\in\N}c_i\le 0,
\]
where $c_0=\P(X_i=0,\,i\in\N)$ and $c_i=\P(X_i>0;\,X_j=0,\,j\ne i)$. For
$m\ge3$, both coefficients $m-1\ge2$ and $m-2\ge1$ are strictly positive,
forcing $c_0=0$ and $c_i=0$ for every $i\in\N$; mass concentrates entirely
on $\M^{(m)}$ with $\sum_{\I\in\M^{(m)}}p_\I=1$. For $m=2$ (ME), the
coefficient of $\sum_i c_i$ vanishes and only $c_0=0$ is implied, with axis
masses remaining unrestricted — consistent with the ME structure. The G
$m$-exclusive case follows identically via Corollary~\ref{cor:m_exclusive_G_star},
with $\F_i(0)$ replaced by $G_i^*$ throughout.

\textit{Parametric families.}
Proposition~\ref{prop:JE_trivariate} provides a one-parameter family for
the trivariate JE ($n=m=3$). For $n=3$ the hierarchy collapses to just ME
($m=2$) and JE ($m=3$), with no non-trivial intermediate level. Non-trivial
intermediate levels first arise at $n=4$ with $m=3$: here $|\M^{(3)}|=6$
unknowns are subject to 4 constraints from~\eqref{eq:JE_parameters_1} and 1
from~\eqref{eq:JE_parameters_2}, so the feasible polytope generically has
dimension greater than one and the one-parameter representation of
Proposition~\ref{prop:JE_trivariate} does not extend. The LP feasibility
argument of Theorem~\ref{thm:m-exclusivity} nonetheless guarantees existence
of at least one admissible solution under~\eqref{eq:m-exclusive_Condition},
with full copula freedom within each admissible face.

A natural further refinement allows specific faces of size $m-1$ to be
deactivated while retaining others, yielding a finer hierarchy governed by
both the global condition~\eqref{eq:m-exclusive_Condition} and local
sub-vector conditions.
\end{remark}
%+++++++++++++++++++++++++++++++++++++++++++++++++++
\section{Further Discussion}
\label{sec:further_discussion}

The preceding sections developed a comprehensive theory of JE and its
extensions. We characterised the admissible marginals, derived an explicit
canonical construction, introduced the G-JE generalisation via marginal
distortion functions, and established the $m$-exclusivity hierarchy
interpolating between ME and JE. We now address four complementary aspects
that complete the picture.

\subsection{Threshold and reflected extensions}
\label{subsec:threshold}

The primary variant studied throughout is \emph{JE from below}, defined for
non-negative $X_\N$ by $\prod_{i\in\N}X_i=0$ almost surely. A symmetric
notion applies to non-positive random vectors: $Y_\N$ is \emph{JE from above}
if $\P(Y_i<0,\,i\in\N)=0$, i.e.\ no $n$ components are simultaneously
strictly negative. Setting $X_i=-Y_i\ge0$, $Y_\N$ is JE from above if and
only if $X_\N$ is JE from below, and condition~\eqref{eq:JECondition} becomes
\[
\sum_{i\in\N}\P(Y_i<0)\le n-1.
\]
All results of Sections~\ref{sec:preliminaries}--\ref{sec:main_results}
transfer immediately by this reflection. The same argument applies at every
level of the $m$-exclusivity hierarchy: a non-positive vector $Y_\N$ is
\emph{$m$-exclusive from above} if $\P(Y_i<0,\,i\in I)=0$ for all
$I\subseteq\N$ with $|I|=m$, which via $X_i=-Y_i$ is equivalent to $X_\N$
being $m$-exclusive from below with existence condition
$\sum_{i\in\N}\P(Y_i<0)\le m-1$.

These cases are unified by a threshold formulation. For thresholds
$l_\N=(l_i,\,i\in\N)\in\mathbb{R}^n$, a random vector $Z_\N$ is
\emph{$m$-exclusive above level $l_\N$} if
\[
\P(Z_i>l_i,\;i\in I)=0
\quad\text{for all }I\subseteq\N\text{ with }|I|=m.
\]
The excess vector $W_i=(Z_i-l_i)_+\ge0$ is then $m$-exclusive from below
with marginal tail probability $\F_{W_i}(0)=\P(Z_i>l_i)$, and the
existence condition for $Z_\N$ with prescribed marginals $\{F_{Z_i}\}$ is
\begin{equation}
\label{eq:m-exclusive_threshold}
\sum_{i\in\N}\P(Z_i>l_i)\le m-1.
\end{equation}
All canonical and G-constructions of Section~\ref{sec:main_results} apply
directly to $W_\N$ and transfer to $Z_\N$ via the shift $Z_i=W_i+l_i$.
The case $m=n$ with $l_\N=0_\N$ and $Z_\N\ge0$ recovers JE from below; the
reflected variant (JE from above on $Y_\N\le0$) is handled via $Z_i=-Y_i$
with $l_\N=0_\N$, as described above, rather than directly through the
threshold formulation. Setting $m=2$ and $l_i=F_{Z_i}^{-1}(q_i)$ for
parameters $q_i\in(0,1)$, condition~\eqref{eq:m-exclusive_threshold} reduces
to $\sum_{i\in\N}(1-q_i)\le1$, the existence condition for ME on the excesses
$W_i=(Z_i-l_i)_+$. The corresponding G-variant
(Corollary~\ref{cor:m_exclusive_G_star}) reduces precisely to the existence
condition for the Tail-ME random vector with probability vector
$(q_1,\dots,q_n)$ studied in~\cite{Cheung2017}, as established in
Example~\ref{ex:G-JE_example}-(1).

\subsection{Aggregate risk and stop-loss implications}
\label{subsec:aggregate}

A central motivation for extremal dependence modelling in insurance
mathematics is the derivation of distributional bounds on aggregate losses.
For a random vector $X_\N$ with fixed non-negative marginals $F_1,\dots,F_n$,
the aggregate $S=\sum_{i\in\N}X_i$ drives excess-of-loss premiums and
solvency capital requirements. Its variance decomposes as
\[
\mathrm{Var}(S)
=\sum_{i\in\N}\mathrm{Var}(X_i)
+2\!\sum_{\substack{i,j\in\N\\i<j}}\!\mathrm{Cov}(X_i,X_j),
\]
where $\mathrm{Cov}(X_i,X_j)=\E[X_iX_j]-\E[X_i]\E[X_j]$. Since the marginals
are held fixed, $\E[X_i]$ and $\mathrm{Var}(X_i)$ are invariant across
dependence structures, and the variance of $S$ is determined entirely by the
pairwise covariances. Under ME, $X_iX_j=0$ almost surely for every $i\ne j$,
so $\E_{ME}[X_iX_j]=0$ and
\[
\mathrm{Cov}_{ME}(X_i,X_j)=-\E[X_i]\E[X_j]\le 0,
\]
giving
\[
\mathrm{Var}(S_{ME})
=\sum_{i\in\N}\mathrm{Var}(X_i)
-2\!\sum_{\substack{i,j\in\N\\i<j}}\!\E[X_i]\E[X_j]
\;\le\;\sum_{i\in\N}\mathrm{Var}(X_i).
\]
Under any $m$-exclusive structure with $m\ge3$, the pairwise constraint is
lifted: no two components are required to satisfy $X_iX_j=0$ almost surely,
and since $X_i,X_j\ge0$ one has $\E[X_iX_j]\ge0$ for all $i\ne j$. The
difference in aggregate variances between any $m$-exclusive ($m\ge3$) and
ME distributions with the same marginals is therefore
\begin{equation}
\label{eq:JE_var_bound}
\mathrm{Var}(S)-\mathrm{Var}(S_{ME})
=2\!\sum_{\substack{i,j\in\N\\i<j}}\!
\bigl(\mathrm{Cov}(X_i,X_j)-\mathrm{Cov}_{ME}(X_i,X_j)\bigr)
=2\!\sum_{\substack{i,j\in\N\\i<j}}\!\E[X_iX_j]\;\ge\;0,
\end{equation}
where the second equality uses $\E_{ME}[X_iX_j]=0$ and the inequality
follows from $X_i,X_j\ge0$. Relaxing ME to any $m$-exclusive structure with
$m\ge3$ therefore weakly increases aggregate variance for fixed marginals,
with equality if and only if $\E[X_iX_j]=0$ for all $i\ne j$, i.e.\ whenever
the structure effectively reduces to ME. When the ME
condition~\eqref{eq:MEcondition} fails but the JE
condition~\eqref{eq:JECondition} holds, the lower
bound~\eqref{eq:JE_parameters_2} forces
$\sum_{\I\in\M}(|\I|-1)p_\I\ge\sum_{i\in\N}\F_i(0)-1>0$, so at least one
face $E_\I$ with $|\I|\ge2$ carries strictly positive mass $p_\I>0$. For
any $i,j\in\I$, $X_i>0$ and $X_j>0$ simultaneously on $E_\I$, and since
$\F_i(0)\ge p_\I>0$ implies $\E[X_i]>0$, we obtain
\[
\E[X_iX_j]
\;\ge\;\E\!\left[X_iX_j\,\mathbf{1}_{E_\I}\right]
=p_\I\,\E\!\left[X_iX_j\,\big|\,E_\I\right]>0,
\]
where the strict inequality holds because $X_i>0$ and $X_j>0$ almost surely
on $E_\I$ and both conditional means are strictly positive. Hence the
inequality in~\eqref{eq:JE_var_bound} is strict for any JE distribution
whose marginals violate the ME condition.

The $m$-exclusivity level also governs the higher-order joint structure of
$S$. Under $m$-exclusivity, $\prod_{i\in I}X_i=0$ almost surely for every
$I\subseteq\N$ with $|I|=m$; since $X_i\ge0$, this implies
$\E[\prod_{i\in I}X_i^{r_i}]=0$ for any $r_i>0$ and all $|I|\ge m$.
Consequently, $m$-exclusivity eliminates all joint moments of order $m$ and
above, directly constraining the cumulants of $S$ of corresponding order.
Under $3$-exclusivity, all triple cross-moments $\E[X_iX_jX_k]$ vanish,
governing the third cumulant of $S$ in addition to the variance; under
$4$-exclusivity, all quartic joint cross-moments vanish as well. This
systematic suppression of higher-order cross-moments as $m$ decreases from
$n$ to $2$ provides a principled and tunable mechanism for controlling the
tail behaviour of the aggregate.

For excess-of-loss reinsurance, the stop-loss premium
$\Pi(d)=\E[(S-d)_+]$ depends on the full distribution of $S$ rather than
its moments alone. When condition~\eqref{eq:MEcondition} holds,
\cite{Cheung2014} establishes that ME minimises $\Pi(d)$ for every retention
level $d\ge0$ among all dependence structures with the given marginals, in
the sense of the convex stochastic order $S_{ME}\le_{cx}S$. This positions
ME as the natural lower benchmark for stop-loss calculations. When the ME
condition fails but the JE condition~\eqref{eq:JECondition} holds, ME is
unavailable as a reference structure. A natural candidate for the stop-loss
minimiser within the JE class is the distribution that suppresses face masses
as far as constraints~\eqref{eq:JE_parameters_1} and \eqref{eq:JE_parameters_2}
permit, minimising $\E[X_iX_j]$ for each pair and approaching ME as closely
as the marginals allow. Whether this candidate achieves a convex ordering
within the JE class, and whether the $m$-exclusivity level $m$ induces a
monotone ordering of stop-loss premiums across the hierarchy for fixed
marginals, are natural and actuarially important open problems.

\subsection{G-JE and distortion operators}
\label{subsec:distortion}

The G-JE construction has a natural interpretation in terms of distortion
operators, a concept central to actuarial risk theory~\citep{Wang1996}. For
each $i\in\N$, the composition $h_i=G_i\circ\F_i$ defines a valid
survival function on $[0,\infty)$: it satisfies $h_i(0)=G_i(\F_i(0))=1$
by the boundary condition $b_i\le\F_i(0)$; it converges to zero as
$x_i\to\infty$ since $\F_i(\infty)=0\le a_i$ and $G_i(0)=0$; and it is
right-continuous and non-increasing. Moreover, since $a_i\ge0$ implies
$G_i(0)=0$, and since $b_i\le\F_i(0)\le1$ implies $G_i(1)=1$, the
function $G_i:[0,1]\to[0,1]$ satisfies the boundary conditions
$G_i(0)=0$ and $G_i(1)=1$ of a distortion function in the sense
of~\cite{Wang1996}, monotonically transforming the survival probabilities of
$X_i$.

Under this identification, a G-JE random vector is a canonical JE random
vector in which the standardised within-face survival function
$\F_i(x_i)/\F_i(0)$ of each component is replaced by the distorted
counterpart $h_i(x_i)=G_i(\F_i(x_i))$. The quantity $G_i^*$ represents
the effective tail mass available for face allocation under the distorted
survival function, and the existence condition~\eqref{eq:G-JECondition} is
condition~\eqref{eq:JECondition} restated in terms of $\{G_i^*\}$ rather than
$\{\F_i(0)\}$. The canonical JE is recovered by setting
$G_i(u)=u/\F_i(0)$ on $[0,\F_i(0)]$ and $G_i(u)=1$ for
$u>\F_i(0)$, giving $G_i^*=\F_i(0)$. Concave distortions
($\gamma\in(0,1)$ in Example~\ref{ex:G-JE_example}-(2)) over-weight the
lower tail of the within-face distribution and yield $G_i^*=0$, causing all
face probabilities to vanish; condition~\eqref{eq:G-JECondition} then reduces
to the ME condition~\eqref{eq:MEcondition}, and the G-JE structure degenerates
to ME. Convex distortions ($\gamma>1$) under-weight the lower tail, giving
$G_i^*=\F_i(0)/\gamma<\F_i(0)$ and tightening the existence
condition relative to canonical JE. This connection links the G-JE framework
naturally to Wang-transform premium principles and to the broader literature
on distortion risk measures, suggesting potential embedding within
distortion-based capital allocation models.

\subsection{Connection to joint mixability}
\label{subsec:JM}

Another prominent extension of counter-monotonicity is \emph{joint
mixability} (JM)~\citep{Wang2016}: a random vector $Y_\N$ is JM if
$\sum_{i\in\N}Y_i=c$ almost surely for some constant $c\in\mathbb{R}$,
taken without loss of generality to be zero. Although JE and JM may appear
structurally unrelated, they share the geometric feature that both supports
are $(n-1)$-dimensional subsets of $\mathbb{R}^n$. Formally, define
\[
W_\mathrm{JE}=\Bigl\{x_\N\in[0,\infty)^n:\min_{i\in\N}x_i=0\Bigr\},
\qquad
W_\mathrm{JM}=\Bigl\{y_\N\in\mathbb{R}^n:\textstyle\sum_{i\in\N}y_i=0\Bigr\},
\]
where $W_\mathrm{JE}$ is the piecewise-linear boundary of the non-negative
orthant and $W_\mathrm{JM}$ is the zero-sum hyperplane. The centering map
$\Phi:W_\mathrm{JE}\to W_\mathrm{JM}$, defined by
\[
\Phi(x_\N)_i=x_i-\xb,\quad i\in\N,\qquad
\xb=\tfrac{1}{n}\sum_{j\in\N}x_j,
\]
has inverse $\Phi^{-1}:W_\mathrm{JM}\to W_\mathrm{JE}$ given by
$\Phi^{-1}(y_\N)_i=y_i-\min_{j\in\N}y_j$. That $\Phi^{-1}$ maps into
$W_\mathrm{JE}$ follows since all components $y_i-\min_j y_j\ge0$ and
$\min_i(y_i-\min_j y_j)=0$. Bijectivity is confirmed by direct computation:
for $x_\N\in W_\mathrm{JE}$ (so $\min_k x_k=0$),
\[
\Phi^{-1}(\Phi(x_\N))_i
=(x_i-\xb)-\min_{k\in\N}(x_k-\xb)
=x_i-\xb-(\min_k x_k-\xb)
=x_i,
\]
and for $y_\N\in W_\mathrm{JM}$ (so $\sum_k y_k=0$),
\[
\Phi(\Phi^{-1}(y_\N))_i
=(y_i-\min_j y_j)
-\tfrac{1}{n}\sum_{k\in\N}(y_k-\min_j y_j)
=y_i-\tfrac{1}{n}\sum_k y_k
=y_i.
\]
Since $\Phi$ is affine and $\Phi^{-1}$ is continuous (as a composition of
continuous functions, the min being continuous), $\Phi$ is a homeomorphism
between $W_\mathrm{JE}$ and $W_\mathrm{JM}$.

The homeomorphism induces a distributional correspondence via pushforward:
if $X_\N$ is JE then $Y_\N:=\Phi(X_\N)=(X_i-\Xb,\,i\in\N)$ is JM,
since $\sum_iY_i=\sum_iX_i-n\Xb=0$ almost surely; and if $Y_\N$ is JM
then $X_\N:=\Phi^{-1}(Y_\N)=(Y_i-\min_{j\in\N}Y_j,\,i\in\N)$ is JE, since
$\min_{i\in\N}X_i=0$ almost surely. An important caveat is that $\Phi$ does
not preserve marginals: the distribution of $Y_i=X_i-\Xb$ depends on
the full joint law of $X_\N$ rather than solely on $F_i$, so the
homeomorphism does not induce a correspondence between JE distributions with
prescribed marginals $\{F_i\}$ and JM distributions with a prescribed
marginal structure. This observation highlights a key analytical advantage
of JE over JM: JE existence is governed by the single tractable
condition~\eqref{eq:JECondition}, whereas JM existence for prescribed
marginals requires substantially more complex analytical conditions;
see~\cite{Wang2016} for a detailed treatment.

For canonical JE random vectors, the homeomorphism carries additional
algebraic structure. On face $E_\I$, where $X_j=0$ for $j\in\N\setminus\I$
and $X_i>0$ for $i\in\I$, the image under $\Phi$ satisfies
\[
\Phi(x_\N)_j=-\tfrac{1}{n}\sum_{k\in\I}x_k
\quad\text{for all }j\in\N\setminus\I,
\]
so all components indexed by $\N\setminus\I$ are equal on $\Phi(E_\I)$. This
block-equality structure in the JM image of each canonical JE face facilitates
explicit computation of JM distributions arising from the canonical
construction and suggests that the bijection may be exploited to transfer
distributional results between the two frameworks. Figure~\ref{fig:JEvsJM}
illustrates $W_\mathrm{JE}$, $W_\mathrm{JM}$, and the action of $\Phi$
for $n=3$.

\begin{figure}[htbp]
\centering
\includegraphics[width=0.75\textwidth]{JEvsJM.pdf}
\caption{The support sets $W_\mathrm{JE}$ and $W_\mathrm{JM}$ for JE and
JM random vectors ($n=3$), and the action of the homeomorphism $\Phi$.}
\label{fig:JEvsJM}
\end{figure}
%+++++++++++++++++++++++++++++++++++++++++++++++++++
\section{Conclusions}
\label{sec:conclusions}

This paper introduces joint exclusivity (JE) as a tractable framework for
extremal negative dependence in arbitrary dimensions. By replacing the
pairwise constraint of mutual exclusivity with the global requirement that
not all $n$ components are simultaneously positive, JE yields a strictly
richer class of dependence structures while retaining a complete and
explicit characterisation: a JE random vector with prescribed marginals
$F_1,\dots,F_n$ exists if and only if $\sum_{i\in\N}\F_i(0)\le n-1$, a
condition whose gap from the mutual exclusivity bound $\sum_{i\in\N}\F_i(0)\le 1$
widens with dimension. The canonical construction separates the allocation
of probability mass across the faces of the support---governed by a
finite-dimensional linear system in the face-mass coefficients---from the
specification of within-face dependence through freely chosen copulas.
In the trivariate case this system reduces to a one-parameter family,
yielding a particularly transparent and implementable model.

Two extensions enrich the framework further. The generalised JE (G-JE)
replaces the standardised within-face survival functions by distorted
counterparts through Wang-type distortion operators $G_i$, accommodating
restricted or reshaped face supports and connecting JE to distortion-based
premium principles. The $m$-exclusivity hierarchy, parametrised by
$2\le m\le n$, interpolates between mutual exclusivity ($m=2$)
and JE ($m=n$), each level characterised by the single inequality
$\sum_{i\in\N}\F_i(0)\le m-1$ and admitting a corresponding G-variant.
Together, these extensions provide a unified and flexible toolkit for
specifying negative dependence across a wide range of modelling contexts.

Several directions emerge naturally for future work. Chief among these is
the stop-loss ordering of JE aggregates: while the variance of $\sum_{i\in\N}X_i$
under JE exceeds that under ME for fixed marginals, a complete characterisation
of excess-of-loss premiums within the JE class remains open and would be of
direct relevance to reinsurance pricing. Equally important are statistical
questions, including the estimation of face-mass coefficients from data and
the development of goodness-of-fit procedures for JE and G-JE models.
Finally, the homeomorphism between the JE and joint mixability supports
suggests that distributional results in each framework may be transferred to
the other, and exploiting this correspondence in the G-JE and $m$-exclusive
settings is a promising avenue for further unifying the theory of extremal
negative dependence.
%%%%%%%%%%%%%%%%%%%%%%%%%%%%%%%%%%%%%%%%%%%%%%%%%%%%%%
\bibliographystyle{apalike}
\bibliography{References.bib}
%\printbibliography
%\bibliographystyle{apalike}
%\bibliography{References2}
%%%%%%%%%%%%%%%%%%%%%%%%%%%%%%%%%%%%%%%%%%%%%%%%%%%%%%%%%%%%%%%%%%%%%%%%%%%%%%%%%%%%%%%%%%%%%%%%%%%%%
\end{document}
